\documentclass[11pt,a4paper]{article}

\usepackage[a4paper,margin=25mm,headheight=14pt]{geometry}
\usepackage{lmodern}
\usepackage{amsmath,amssymb,amsthm,mathtools}
\usepackage{booktabs}
\usepackage{tabularx}
\usepackage{fancyhdr}
\usepackage{xcolor}
\usepackage{microtype}
\usepackage{hyperref}

\definecolor{primary}{RGB}{20,62,115}
\hypersetup{
  colorlinks=true,
  linkcolor=primary,
  citecolor=primary,
  urlcolor=primary,
  pdftitle={A Spectral Dichotomy for Dilworth Multiplicity Bounds of Q(G) and D^Q(G) on C-Graphs and Cographs}
}

\setlength{\parskip}{5pt plus 1.5pt minus 1pt}
\setlength{\parindent}{1.5em}

\pagestyle{fancy}
\fancyhf{}
\fancyhead[L]{\small\itshape Dilworth Multiplicity Dichotomy for $Q(G)$ and $D^Q(G)$ on $\mathcal{C}$-Graphs and Cographs}
\fancyhead[R]{\small\thepage}
\renewcommand{\headrulewidth}{0.4pt}

\newtheorem{theorem}{Theorem}[section]
\newtheorem{corollary}[theorem]{Corollary}
\newtheorem{proposition}[theorem]{Proposition}
\newtheorem{lemma}[theorem]{Lemma}
\newtheorem{conjecture}[theorem]{Conjecture}
\newtheorem{problem}[theorem]{Open Problem}
\theoremstyle{definition}
\newtheorem{definition}[theorem]{Definition}
\newtheorem{example}[theorem]{Example}
\newtheorem{remark}[theorem]{Remark}

\DeclareMathOperator{\Tr}{Tr}
\DeclareMathOperator{\mult}{mult}
\DeclareMathOperator{\spec}{spec}
\DeclareMathOperator{\Deg}{Deg}
\DeclareMathOperator{\diam}{diam}
\DeclareMathOperator{\rank}{rank}
\DeclareMathOperator{\nullity}{nullity}
\DeclareMathOperator{\im}{im}

\title{\Large\textbf{A Spectral Dichotomy for Dilworth Multiplicity Bounds of the Signless Laplacian and Distance Signless Laplacian on $\mathcal{C}$-Graphs and Cographs}}
\author{Research Manuscript}
\date{November 27, 2026}

\begin{document}
\maketitle
\thispagestyle{plain}

\begin{abstract}
\noindent
For the adjacency matrix $A(G)$, distance matrix $D(G)$, and Seidel matrix $S(G)$ of any connected cograph $G$, every non-trivial eigenvalue has multiplicity at most the Dilworth number $\nabla(G)$ (Ghorbani~2019; Lou--Lin~2021; Li--Wang--Brunetti~2024). Motivated by this unity, Mandal and Manna (\href{https://arxiv.org/abs/2610.08599}{arXiv:2610.08599}, 2026; reiterated as Open Problem~14.6 in \href{https://arxiv.org/abs/2611.06981}{arXiv:2611.06981}) conjectured that every block-non-local eigenvalue $\theta \notin \mathcal{T}^Q(G)$ of the distance signless Laplacian $D^Q(G)$ of a connected $\mathcal{C}$-graph satisfies $\mult_{D^Q(G)}(\theta) \le \nabla(G)$. We establish a sharp structural dichotomy between $\mathcal{C}$-graphs and general cographs for both the signless Laplacian $Q(G)$ and the distance signless Laplacian $D^Q(G)$:
\begin{enumerate}
  \item[(i)] \emph{Exact Dilworth number and refutation on $\mathcal{C}$-graphs:} For every connected $\mathcal{C}$-graph $G = C(\alpha_1,\dots,\alpha_m)$ ($m \ge 2$ even), $\nabla(G) = 1 + |\{j \ge 3 : \alpha_j \ge 2\}|$. We construct an infinite counterexample family $G_k = C(k,1,\dots,1)$ ($n = 3k+5$, $k \ge 1$) with $\nabla(G_k)=1$ and $\mult_{D^Q(G_k)}(4k+4) = 2$ (complementing a prior independent Lean~4 refutation of Conjecture~5.5 recorded in \texttt{trureturing}), and verify via exact symbolic irreducible-factor decomposition in SymPy over all $127$ even-length $\mathcal{C}$-graph creation sequences of order $n \le 8$ that the $8$-vertex antiregular graph $H_8 = C(1^8)$ is the unique smallest counterexample up to isomorphism.
  \item[(ii)] \emph{Universal sharp bounds for all $\mathcal{C}$-graphs:} Via an exact block recurrence, an explicit free-parameter isomorphism with the quotient eigenspace, critical-value separation, and a formal conservation-law tail propagation argument, we prove analytically that \emph{every} connected $\mathcal{C}$-graph satisfies $\mult_{Q(G)}(\mu) \le \nabla(G)$ for all $\mu \notin \mathcal{T}_Q(G)$ and, via a unified rank-one perturbation lemma, the sharp repaired bound $\mult_{D^Q(G)}(\theta) \le \nabla(G) + 1$ for all $\theta \notin \mathcal{T}^Q(G)$.
  \item[(iii)] \emph{General cographs—universal linear upper bounds and linear separation:} For maximal-twin non-local eigenvalues ($\mu \notin \mathcal{T}_{Q,\max}(G)$, $\theta \notin \mathcal{T}^Q_{\max}(G)$), we prove that every non-local eigenvalue of an integer vicinal-monotone vertex-weighted signless Laplacian on any connected threshold graph is simple. Combining a canonical cotree shell--core decomposition lemma with weighted twin and local eigenvalue inheritance lemmas, we prove analytically that \emph{every} connected cograph $G$ satisfies
  \begin{align*}
    \mult_{Q(G)}(\mu) &\le 5\nabla(G) - 4 && (\mu \notin \mathcal{T}_{Q,\max}(G)), \\
    \mult_{D^Q(G)}(\theta) &\le 5\nabla(G) - 3 && (\theta \notin \mathcal{T}^Q_{\max}(G)).
  \end{align*}
  Conversely, although exact symbolic enumeration in SymPy of all $1{,}170$ non-isomorphic connected cographs of order $n \le 9$ reveals zero violations of $\mult_Q \le \nabla$ and $\mult_{D^Q} \le \nabla + 1$, \emph{both} bounds fail linearly on general connected cographs: starting from the $60$-vertex cograph $B = K_{36} \vee (3 H_8)$ (for which we prove analytically that $\nabla(B)=3$ and $\mult_{Q(B)}(40) = 4$), the family $F_s = K_1 \vee (s \cdot B)$ ($s \ge 1$) satisfies $\nabla(F_s) = 3s$ while both $Q(F_s)$ and $D^Q(F_s)$ attain non-local multiplicity at least $4s - 1 = \frac{4}{3}\nabla(F_s) - 1$. Defining $M_Q(k)$ and $M_{D^Q}(k)$ as the maximum non-local multiplicities over connected cographs of Dilworth number $k$, the asymptotic growth constants satisfy $\frac{4}{3} \le c_Q^* = c_{D^Q}^* \le 5$, separating $Q$ and $D^Q$ fundamentally from $A$, $D$, and $S$.
\end{enumerate}

\smallskip
\noindent\textbf{Keywords:} Distance signless Laplacian; signless Laplacian; $\mathcal{C}$-graph; cograph; threshold graph; Dilworth number; eigenvalue multiplicity.

\smallskip
\noindent\textbf{2020 Mathematics Subject Classification:} 05C50; 15A18.
\end{abstract}

\section{Introduction and Preliminaries}\label{sec:intro}

Let $G = (V, E)$ be a finite, simple undirected graph of order $n = |V|$ with adjacency matrix $A(G)$, diagonal degree matrix $\Deg(G)$, and signless Laplacian $Q(G) = \Deg(G) + A(G)$. When $G$ is connected, let $D(G) = (d(u,v))_{u,v\in V}$ be its distance matrix, $\Tr(G) = \mathrm{diag}(\Tr(v))_{v\in V}$ its diagonal transmission matrix ($\Tr(v) = \sum_{u\in V} d(u,v)$), and $D^Q(G) = \Tr(G) + D(G)$ its \emph{distance signless Laplacian}~\cite{AH2016}. For any connected graph with diameter $\diam(G) \le 2$ (in particular, every connected cograph), $D(G) = 2(J-I) - A(G)$ and $\Tr(G) = 2(n-1)I - \Deg(G)$, yielding the rank-one perturbation identity
\begin{equation}\label{eq:rank1}
  D^Q(G) = (2n - 4)I - Q(G) + 2J,
\end{equation}
where $J = \mathbf{1}\mathbf{1}^T$ is the all-ones matrix~\cite{MM2026,MMT2026}.

A unifying paradigm in the spectral theory of cographs ($P_4$-free graphs) is bounding non-trivial eigenvalue multiplicities by the \emph{Dilworth number} $\nabla(G)$. Under the \emph{vicinal preorder} $u \preceq v \iff N(u) \subseteq N[v]$, $\nabla(G)$ is the minimum number of chains covering $V(G)$ (equivalently, by Dilworth's theorem on the induced poset of twin classes, the maximum size of an antichain); a graph satisfies $\nabla(G) = 1$ if and only if it is a \emph{threshold graph} ($\{P_4, C_4, 2K_2\}$-free). Across three major graph matrices, the exact bound $\mult \le \nabla(G)$ holds universally for all connected cographs: for the adjacency matrix $A(G)$ ($\lambda \notin \{0,-1\}$, Ghorbani~\cite{Ghorbani2019}), the distance matrix $D(G)$ ($\lambda \notin \{-2,-1\}$, Lou and Lin~\cite{LL2021}), and the Seidel matrix $S(G)$ ($\lambda \notin \{-1,1\}$, Li, Wang, and Brunetti~\cite{LWB2024}).

\subsection{Twin Classes, Local Eigenvalue Sets, and Equitable Quotients}

Two distinct vertices $u, v \in V(G)$ are \emph{twins} if $N(u) \setminus \{v\} = N(v) \setminus \{u\}$ (\emph{true/clique twins} if $uv \in E$, \emph{false/independent twins} if $uv \notin E$). Being equal or twins is an equivalence relation whose equivalence classes are the \emph{maximal twin classes} of $G$. Each maximal twin class $C$ of size $|C| \ge 2$ is either a clique or an independent set of common degree $d_C$ and transmission $\Tr_C = 2(n-1)-d_C$ (when $\diam(G) \le 2$), contributing $|C|-1$ linearly independent twin-difference eigenvectors supported on $C$ with \emph{maximal-twin-local eigenvalues}:
\begin{align*}
  \mathcal{T}_{Q,\max}(G) &:= \{d_C - 1 : C\text{ clique},\, |C| \ge 2\} \cup \{d_C : C\text{ independent},\, |C| \ge 2\}, \\
  \mathcal{T}^Q_{\max}(G) &:= \{\Tr_C - 1 : C\text{ clique},\, |C| \ge 2\} \cup \{\Tr_C - 2 : C\text{ independent},\, |C| \ge 2\}.
\end{align*}

Mandal and Mehatari~\cite{MM2025} introduced the subclass of \emph{$\mathcal{C}$-graphs}, generated from a finite creation sequence $(\alpha_1, \dots, \alpha_m)$ of positive integers summing to $n$ by
\[
  C(\alpha_1) = \overline{K_{\alpha_1}}, \qquad
  C(\alpha_1,\dots,\alpha_i) = \overline{\,C(\alpha_1,\dots,\alpha_{i-1}) \cup K_{\alpha_i}\,} \quad (i \ge 2).
\]
Every $\mathcal{C}$-graph with $m \ge 2$ even is connected, has $\diam(G) \le 2$, and partitions $V(G)$ into blocks $C_1,\dots,C_m$ ($|C_i| = \alpha_i$) such that~\cite[Proposition~4.2]{MM2026}:
\begin{itemize}
  \item $C_i$ is a clique if $i$ is odd, and an independent set if $i$ is even;
  \item for $1 \le i < j \le m$, $C_i$ is completely adjacent to $C_j$ if $j$ is even, and completely non-adjacent to $C_j$ if $j$ is odd.
\end{itemize}
Following~\cite{MM2026}, blocks $C_i$ with $\alpha_i \ge 2$ define the \emph{block-local eigenvalue sets}
\begin{align*}
  \mathcal{T}_Q(G) &:= \{d_i - 1 : i\text{ odd},\, \alpha_i \ge 2\} \cup \{d_i : i\text{ even},\, \alpha_i \ge 2\}, \\
  \mathcal{T}^Q(G) &:= \{\Tr_i - 1 : i\text{ odd},\, \alpha_i \ge 2\} \cup \{\Tr_i - 2 : i\text{ even},\, \alpha_i \ge 2\}.
\end{align*}
The remaining $m$ eigenvalues of $Q(G)$ and $D^Q(G)$ are those of the $m \times m$ equitable quotient matrices $Q_\pi(G)$ and $Q_{D^Q}(G)$~\cite[Theorem~4.3]{MM2026}:
\begin{equation}\label{eq:Qpi}
  Q_\pi(i,i) = d_i + (\alpha_i - 1)\delta_i, \qquad
  Q_\pi(i,j) = \alpha_j \delta_{ij} \quad (i \ne j),
\end{equation}
where $\delta_i = 1$ ($i$ odd) or $0$ ($i$ even), and $\delta_{ij} = 1$ ($\max(i,j)$ even) or $0$ ($\max(i,j)$ odd); and
\begin{equation}\label{eq:QDQ}
  Q_{D^Q}(i,i) = \Tr_i + (\alpha_i - 1) c_i, \qquad
  Q_{D^Q}(i,j) = \alpha_j c_{ij} \quad (i \ne j),
\end{equation}
where $c_i = 1$ ($i$ odd) or $2$ ($i$ even), and $c_{ij} = 1$ ($\max(i,j)$ even) or $2$ ($\max(i,j)$ odd).

\begin{lemma}[Diagonalizability of Equitable Quotient Matrices]\label{lem:diag}
  For any connected $\mathcal{C}$-graph $G = C(\alpha_1,\dots,\alpha_m)$ ($m \ge 2$ even), let $R := \mathrm{diag}(\sqrt{\alpha_1},\dots,\sqrt{\alpha_m})$. Then $R Q_\pi(G) R^{-1}$ and $R Q_{D^Q}(G) R^{-1}$ are real symmetric matrices. In particular, $Q_\pi(G)$ and $Q_{D^Q}(G)$ are diagonalizable over $\mathbb{R}$, so for every $\mu, \theta \in \mathbb{R}$,
  \[
    \mult_{Q_\pi(G)}(\mu) = \dim\ker(Q_\pi(G) - \mu I_m), \qquad
    \mult_{Q_{D^Q}(G)}(\theta) = \dim\ker(Q_{D^Q}(G) - \theta I_m).
  \]
\end{lemma}

\begin{proof}
  For $i \ne j$, $(R Q_\pi R^{-1})_{ij} = \sqrt{\alpha_i / \alpha_j}\,\alpha_j \delta_{ij} = \sqrt{\alpha_i \alpha_j}\,\delta_{ij}$, which is symmetric in $i, j$ because $\delta_{ij}$ depends only on $\max(i,j)$. Identically, $(R Q_{D^Q} R^{-1})_{ij} = \sqrt{\alpha_i \alpha_j}\,c_{ij}$ is symmetric in $i, j$. Since every real symmetric matrix is orthogonally diagonalizable, $Q_\pi$ and $Q_{D^Q}$ are diagonalizable over $\mathbb{R}$, so algebraic and geometric multiplicities coincide.
\end{proof}

\begin{proposition}[Relation Between Block-Local and Maximal-Twin-Local Sets]\label{prop:local_rel}
  Let $G = C(\alpha_1,\dots,\alpha_m)$ be a connected $\mathcal{C}$-graph ($m \ge 2$ even).
  \begin{enumerate}
    \item The blocks $C_3,\dots,C_m$ are pairwise distinct maximal twin classes of $G$. Moreover, $C_1$ and $C_2$ merge into a single maximal clique twin class $C_1 \cup C_2$ if and only if $\alpha_2 = 1$; otherwise $C_1$ and $C_2$ are distinct maximal twin classes.
    \item Consequently, $\mathcal{T}_Q(G) \subseteq \mathcal{T}_{Q,\max}(G)$ and $\mathcal{T}^Q(G) \subseteq \mathcal{T}^Q_{\max}(G)$, with equality in both inclusions whenever $\alpha_2 \ge 2$ or $\alpha_1 \ge 2$. When $\alpha_1 = \alpha_2 = 1$,
    \[
      \mathcal{T}_{Q,\max}(G) = \mathcal{T}_Q(G) \cup \{d_1 - 1\}, \qquad
      \mathcal{T}^Q_{\max}(G) = \mathcal{T}^Q(G) \cup \{\Tr_1 - 1\}.
    \]
  \end{enumerate}
\end{proposition}

\begin{proof}
  Vertices in $C_1$ and $C_2$ have identical adjacencies to all $C_r$ ($r \ge 3$), and $C_1 \sim C_2$. When $\alpha_2 = 1$, $C_1 \cup C_2$ is a clique whose vertices share the exact same closed neighborhood, so $C_1 \cup C_2$ is a clique twin class (of degree $d_1 = d_2$, contributing $d_1 - 1$ to $\mathcal{T}_{Q,\max}$ and $\Tr_1 - 1$ to $\mathcal{T}^Q_{\max}$); when $\alpha_2 \ge 2$, vertices in $C_2$ are non-adjacent whereas vertices in $C_1$ are adjacent to all of $C_2$, so $C_1$ and $C_2$ are distinct maximal twin classes. For any $j \ge 3$ and $i < j$, if $j$ is odd then $C_j$ is non-adjacent to $C_1,\dots,C_{j-1}$ while $C_i$ has at least one neighbor in $C_1 \cup \dots \cup C_{j-1}$ (namely $C_2$ if $i=1$, $C_1$ if $i \ge 2$ is even, or $C_{i+1}$ if $i \ge 3$ is odd); if $j$ is even then $C_j$ is adjacent to $C_{j-1}$ while $C_i$ is non-adjacent to $C_{j-1}$ (if $i < j-1$), or $C_j$ is adjacent to $C_1$ while $C_{j-1}$ is non-adjacent to $C_1$ (if $i = j-1 \ge 3$). Hence no block $C_j$ ($j \ge 3$) merges with any other block. Part~(2) follows directly: when $\alpha_1 \ge 2$ and $\alpha_2 = 1$, $d_1 - 1$ is already in $\mathcal{T}_Q(G)$, whereas when $\alpha_1 = \alpha_2 = 1$, $d_1 - 1$ lies in $\mathcal{T}_{Q,\max}(G) \setminus \mathcal{T}_Q(G)$ unless some other block happens to share the same local eigenvalue.
\end{proof}

Throughout this paper, we state explicitly which local eigenvalue set applies: Section~\ref{sec:Cgraph_bound} establishes bounds for the stronger \emph{block-non-local} conditions $\mu \notin \mathcal{T}_Q(G)$ and $\theta \notin \mathcal{T}^Q(G)$ (matching Conjecture~\ref{conj:MM} verbatim, and a fortiori implying the maximal-twin-non-local versions by Proposition~\ref{prop:local_rel}), whereas Section~\ref{sec:cographs} works uniformly with the \emph{maximal-twin-non-local} sets $\mathcal{T}_{Q,\max}(G)$ and $\mathcal{T}^Q_{\max}(G)$ on general cographs.

\begin{lemma}[Unified Rank-One Perturbation and Local Eigenvalue Correspondence]\label{lem:rank1}
  \mbox{}
  \begin{enumerate}
    \item For any real $n \times n$ matrices $M, N$ with $\rank(M - N) \le r$,
    \[
      \bigl|\nullity(M) - \nullity(N)\bigr| \le r.
    \]
    \item For every connected graph $G$ of order $n$ with $\diam(G) \le 2$ and every $\theta \in \mathbb{R}$,
    \[
      \bigl|\mult_{D^Q(G)}(\theta) - \mult_{Q(G)}(2n - 4 - \theta)\bigr| \le 1.
    \]
    \item Under the involution $\theta = 2n - 4 - \mu$, for every connected graph $G$ with $\diam(G) \le 2$,
    \[
      \theta \notin \mathcal{T}^Q_{\max}(G) \iff 2n - 4 - \theta \notin \mathcal{T}_{Q,\max}(G),
    \]
    and for every connected $\mathcal{C}$-graph $G$, $\theta \notin \mathcal{T}^Q(G) \iff 2n - 4 - \theta \notin \mathcal{T}_Q(G)$.
  \end{enumerate}
\end{lemma}

\begin{proof}
  \textbf{(1)} By subadditivity of rank, $\rank(M) \le \rank(N) + \rank(M-N) \le \rank(N) + r$ and symmetrically $\rank(N) \le \rank(M) + r$, so $|\rank(M) - \rank(N)| \le r$; rank-nullity ($\nullity = n - \rank$) gives $|\nullity(M) - \nullity(N)| \le r$.

  \textbf{(2)} Both $Q(G)$ and $D^Q(G)$ are real symmetric, so algebraic and geometric multiplicities coincide. By~\eqref{eq:rank1}, $(D^Q(G) - \theta I) - ((2n-4-\theta)I - Q(G)) = 2J$ has rank~$1$, so~(1) yields $|\mult_{D^Q(G)}(\theta) - \mult_{Q(G)}(2n-4-\theta)| \le 1$.

  \textbf{(3)} Since $\diam(G) \le 2$, every twin class $C$ of common degree $d_C$ has transmission $\Tr_C = 2(n-1) - d_C$. For a clique twin class, $2n - 4 - (d_C - 1) = 2n - 3 - d_C = \Tr_C - 1$; for an independent twin class, $2n - 4 - d_C = \Tr_C - 2$. Thus $\mu \mapsto 2n - 4 - \mu$ maps $\mathcal{T}_{Q,\max}(G)$ bijectively onto $\mathcal{T}^Q_{\max}(G)$ (and $\mathcal{T}_Q(G)$ bijectively onto $\mathcal{T}^Q(G)$ for $\mathcal{C}$-graphs).
\end{proof}

\subsection{Target Conjecture, Prior Work, and Summary of Results}

Mandal and Manna~\cite[Conjecture~5.5]{MM2026} (reiterated as Conjecture~10.10 and Open Problem~14.6 in~\cite{MMT2026}) conjectured:

\begin{conjecture}[Mandal--Manna~\cite{MM2026}, Conjecture~5.5]\label{conj:MM}
  For every connected $\mathcal{C}$-graph $G$, every eigenvalue $\theta$ of $D^Q(G)$ not in $\mathcal{T}^Q(G)$ satisfies $\mult_{D^Q(G)}(\theta) \le \nabla(G)$.
\end{conjecture}

\paragraph{Relation to prior work and scope of verification.}
An independent Lean~4 refutation of Conjecture~\ref{conj:MM} was recorded on November~7, 2026 in the \texttt{trureturing} external open problems registry~\cite{TrureTuring2026}. In the present paper, all general structural theorems (Theorems~\ref{thm:dilworth}, \ref{thm:disproof}, \ref{thm:CgraphQ}, \ref{thm:CgraphDQ}, \ref{thm:Qsimple}, \ref{thm:cograph_upper}, and~\ref{thm:cograph_lower}) are proved purely analytically. In addition, finite-order enumerations (Proposition~\ref{prop:minimality} and Table~\ref{tab:cographs}) and sanity checks of our explicit constructions are performed via a companion pure symbolic Python/SymPy script (\texttt{verify\_all.py}) with zero floating-point arithmetic; we emphasize that finite-order symbolic checks confirm exact small-order behavior and concrete instances, whereas all universal bounds rest entirely on our analytical proofs.
\begin{itemize}
  \item \textbf{Section~\ref{sec:dilworth}:} Exact Dilworth number $\nabla(C(\alpha_1,\dots,\alpha_m)) = 1 + |\{j \ge 3 : \alpha_j \ge 2\}|$ (Theorem~\ref{thm:dilworth}), refining Remark~2.4 of~\cite{MMT2026}.
  \item \textbf{Section~\ref{sec:disproof}:} Infinite counterexample family $G_k = C(k,1,\dots,1)$ ($k \ge 1$) with $\nabla(G_k)=1$ and $\mult_{D^Q(G_k)}(4k+4) = 2$ (Theorem~\ref{thm:disproof}), plus exact symbolic proof over all $127$ even-length $\mathcal{C}$-graph creation sequences of order $n \le 8$ that $H_8 = C(1^8)$ is the unique smallest counterexample up to isomorphism (Proposition~\ref{prop:minimality}).
  \item \textbf{Section~\ref{sec:Cgraph_bound}:} Complete analytical proof that \emph{every} connected $\mathcal{C}$-graph satisfies $\mult_{Q(G)}(\mu) \le \nabla(G)$ for $\mu \notin \mathcal{T}_Q(G)$ (Theorem~\ref{thm:CgraphQ}) and the sharp repaired bound $\mult_{D^Q(G)}(\theta) \le \nabla(G) + 1$ for $\theta \notin \mathcal{T}^Q(G)$ (Theorem~\ref{thm:CgraphDQ}).
  \item \textbf{Section~\ref{sec:cographs}:} Unified cotree formula for $\nabla(G)$ (Lemma~\ref{lem:cotree_nabla}), canonical block form and weighted simplicity for connected threshold graphs (Lemma~\ref{lem:thresh_canon}, Theorem~\ref{thm:Qsimple}), standalone canonical cotree shell--core decomposition, twin-inheritance, and shell-restriction lemmas (Lemmas~\ref{lem:shell_core}--\ref{lem:shell_restrict}), self-contained proof of the universal linear upper bounds $\mult_{Q(G)}(\mu) \le 5\nabla(G) - 4$ and $\mult_{D^Q(G)}(\theta) \le 5\nabla(G) - 3$ on all connected cographs (Theorem~\ref{thm:cograph_upper}), and an analytical proof of linear separation $\frac{4}{3}\nabla(F_s) - 1$ on general connected cographs (Theorem~\ref{thm:cograph_lower}), yielding $\frac{4}{3} \le c_Q^* = c_{D^Q}^* \le 5$ (Open Problem~\ref{prob:cstar}).
\end{itemize}

\begin{table}[!htbp]
  \centering
  \caption{Summary of non-local eigenvalue multiplicity bounds by graph matrix and class (block-non-local $\mathcal{T}_Q, \mathcal{T}^Q$ on $\mathcal{C}$-graphs; maximal-twin-non-local $\mathcal{T}_{Q,\max}, \mathcal{T}^Q_{\max}$ on threshold graphs and general cographs).}
  \label{tab:dichotomy}
  \begin{tabularx}{\linewidth}{lXX>{\raggedright\arraybackslash}X}
    \toprule
    Matrix & Threshold ($\nabla=1$) & Connected $\mathcal{C}$-graphs & Connected cographs \\
    \midrule
    $A, D, S$ & $\le 1$ (sharp) & $\le \nabla$ & $\le \nabla$~\cite{Ghorbani2019,LL2021,LWB2024} \\
    $Q(G)$ & $\le 1$ (Thm.~\ref{thm:Qsimple}) & $\le \nabla$ (Thm.~\ref{thm:CgraphQ}, sharp) & $\le 5\nabla - 4$ (Thm.~\ref{thm:cograph_upper});\newline $\ge \frac{4}{3}\nabla - 1$ (Thm.~\ref{thm:cograph_lower}) \\
    $D^Q(G)$ & $\le 2$ (Cor.~\ref{cor:repaired}, sharp) & $\le \nabla + 1$ (Thm.~\ref{thm:CgraphDQ}, sharp) & $\le 5\nabla - 3$ (Thm.~\ref{thm:cograph_upper});\newline $\ge \frac{4}{3}\nabla - 1$ (Thm.~\ref{thm:cograph_lower}) \\
    \bottomrule
  \end{tabularx}
\end{table}

\section{Exact Dilworth Number of $\mathcal{C}$-Graphs}\label{sec:dilworth}

Because twins are mutually comparable in the vicinal preorder, $\nabla(G)$ is determined by the induced preorder on the blocks $C_1,\dots,C_m$.

\begin{theorem}\label{thm:dilworth}
  Let $G = C(\alpha_1,\dots,\alpha_m)$ be a connected $\mathcal{C}$-graph ($m \ge 2$ even). Then
  \[
    \nabla(G) = 1 + \bigl|\{j \in \{3, \dots, m\} : \alpha_j \ge 2\}\bigr|.
  \]
\end{theorem}

\begin{proof}
  Write $C_i \preceq C_j$ if $u \preceq v$ for all $u \in C_i$ and $v \in C_j$. Let $S = \{j \in \{3,\dots,m\} : \alpha_j \ge 2\}$.

  First, since $C_1$ is a clique, $C_2$ is an independent set, $C_1 \sim C_2$, and $C_1, C_2$ share identical adjacencies to every $C_r$ ($r \ge 3$), for any $u \in C_2$ and $v \in C_1$ we have
  \[
    N(u) = C_1 \cup \bigcup_{\substack{r \ge 4 \\ r\text{ even}}} C_r \;\subseteq\; C_1 \cup C_2 \cup \bigcup_{\substack{r \ge 4 \\ r\text{ even}}} C_r = N[v],
  \]
  so $C_2 \preceq C_1$ always holds.

  Next, let $1 \le i < j \le m$ with $j \ge 3$, and pick $u \in C_i$, $v \in C_j$.
  \begin{itemize}
    \item \emph{Case 1: $j$ odd ($j \ge 3$).} Here $C_j$ is a clique non-adjacent to $C_1,\dots,C_{j-1}$ and adjacent to all even blocks $C_r$ with $r > j$. If $\alpha_j = 1$, then $N(v) = \bigcup_{r>j,\, r\text{ even}} C_r \subseteq N(u) \subseteq N[u]$, so $C_j \preceq C_i$. If $\alpha_j \ge 2$, then $v$ has a neighbor $v' \in C_j \setminus \{v\}$ non-adjacent to $u$, so $C_j \not\preceq C_i$; conversely, $u$ is adjacent to $C_2$ (if $i=1$), to $C_1$ (if $i \ge 2$ is even), or to $C_{i+1}$ (if $i \ge 3$ is odd, since $i+1 < j$ is even), whereas $v$ is non-adjacent to all of $C_1,\dots,C_{j-1}$, so $C_i \not\preceq C_j$. Thus $C_i$ and $C_j$ are incomparable.
    \item \emph{Case 2: $j$ even ($j \ge 4$).} Here $C_j$ is an independent set adjacent to $C_1,\dots,C_{j-1}$ and to all even blocks $C_r$ with $r > j$. If $\alpha_j = 1$, then $N[v] = C_1 \cup \dots \cup C_{j-1} \cup \{v\} \cup \bigcup_{r>j,\, r\text{ even}} C_r \supseteq N(u)$, so $C_i \preceq C_j$. If $\alpha_j \ge 2$, then $u$ is adjacent to every $v' \in C_j \setminus \{v\}$ whereas $v$ is not, so $C_i \not\preceq C_j$; conversely, $v$ is adjacent to $C_{j-1}$ (non-adjacent to $u$ when $i < j-1$) and to $C_1$ (non-adjacent to $u$ when $i = j-1 \ge 3$), so $C_j \not\preceq C_i$. Thus $C_i$ and $C_j$ are incomparable.
  \end{itemize}

  \emph{Lower bound.} Selecting one representative $v_1 \in C_1$ and one representative $v_j \in C_j$ for each $j \in S$ yields an antichain of size $1 + |S|$; hence $\nabla(G) \ge 1 + |S|$.

  \emph{Upper bound.} Each $C_j$ ($j \in S$) is a twin class, hence a chain. Writing $O_1^{\mathrm{odd}} = \{j \ge 3 : j\text{ odd},\, \alpha_j = 1\}$ and $O_1^{\mathrm{even}} = \{j \ge 4 : j\text{ even},\, \alpha_j = 1\}$, all remaining vertices $V(G) \setminus \bigcup_{j\in S} C_j$ form a single chain:
  \[
    \bigl(\text{blocks in } O_1^{\mathrm{odd}}\text{ decreasing}\bigr) \;\preceq\; C_2 \;\preceq\; C_1 \;\preceq\; \bigl(\text{blocks in } O_1^{\mathrm{even}}\text{ increasing}\bigr).
  \]
  Thus $V(G)$ is covered by $1 + |S|$ chains, proving $\nabla(G) \le 1 + |S|$.
\end{proof}

\begin{corollary}\label{cor:threshold}
  A connected $\mathcal{C}$-graph $G = C(\alpha_1,\dots,\alpha_m)$ ($m \ge 2$ even) is a threshold graph if and only if $\alpha_3 = \alpha_4 = \dots = \alpha_m = 1$.
\end{corollary}

\begin{remark}
  Corollary~\ref{cor:threshold} refines Remark~2.4 of~\cite{MMT2026}, which stated that threshold graphs and $\mathcal{C}$-graphs meet only in the antiregular graphs $H_{2k} = C(1,\dots,1)$ and $H_{2k+1} = C(1,2,1,\dots,1)$; in fact, every $C(\alpha_1,\alpha_2,1,\dots,1)$ with $\alpha_1,\alpha_2 \ge 1$ is a threshold graph.
\end{remark}

\section{Infinite Counterexample Family and Unique Minimality}\label{sec:disproof}

For each integer $k \ge 1$, set $t = k+3$ and $m = 2t = 2k+6$, and define the connected $\mathcal{C}$-graph
\[
  G_k = C(k, \underbrace{1, 1, \dots, 1}_{2k+5})
\]
of order $n = 3k + 5$. For $k = 1$, $G_1 = C(1^8)$ is the $8$-vertex antiregular graph $H_8$.

\begin{theorem}\label{thm:disproof}
  For every integer $k \ge 1$, the connected $\mathcal{C}$-graph $G_k$ on $n = 3k+5$ vertices satisfies:
  \begin{enumerate}
    \item $\nabla(G_k) = 1$;
    \item $\theta_k := 4k + 4 \notin \mathcal{T}^Q(G_k)$, and $\theta_k \notin \mathcal{T}^Q_{\max}(G_k)$;
    \item $\mult_{D^Q(G_k)}(4k + 4) = 2 > 1 = \nabla(G_k)$.
  \end{enumerate}
  Consequently, Conjecture~\ref{conj:MM} is false for infinitely many connected $\mathcal{C}$-graphs.
\end{theorem}

\begin{proof}
  \textbf{(1)} Immediate from Theorem~\ref{thm:dilworth}.

  \textbf{(2)} With $t = k+3$ and $n = 3k+5$, the block degrees $d_i$ and transmissions $\Tr_i = 6k+8-d_i$ are:
  \begin{align*}
    d_1 &= (k-1) + t = 2k+2, & \Tr_1 &= 4k+6, \\
    d_2 &= k + (t-1) = 2k+2, & \Tr_2 &= 4k+6, \\
    d_{2j-1} &= t - j + 1 = k + 4 - j, & \Tr_{2j-1} &= 5k + 4 + j && (2 \le j \le k+3), \\
    d_{2j} &= (k + 2j - 2) + (t - j) = 2k + 1 + j, & \Tr_{2j} &= 4k + 7 - j && (2 \le j \le k+3).
  \end{align*}
  Thus $\mathcal{T}^Q(G_1) = \emptyset$ and $\mathcal{T}^Q(G_k) = \{4k+5\}$ for $k \ge 2$, while by Proposition~\ref{prop:local_rel} the unique maximal twin class of size $\ge 2$ is $C_1 \cup C_2$ with $\mathcal{T}^Q_{\max}(G_k) = \{4k+5\}$ for all $k \ge 1$. Hence $\theta_k = 4k+4$ is non-local under both definitions.

  \textbf{(3)} Define $\mathbf{x}^{(1)}, \mathbf{x}^{(2)} \in \mathbb{R}^{2k+6}$ by
  \begin{align*}
    \mathbf{x}^{(1)} &= (-1,\,-1,\,1,\,k,\,\underbrace{0,0,\dots,0}_{2k+2})^T, \\
    \mathbf{x}^{(2)} &= \bigl(k+1,\,k+1,\,-(k+1),\,0,\,\underbrace{-k,\,k,\,-k,\,k,\,\dots,\,-k,\,k}_{k+1\text{ pairs}}\bigr)^T,
  \end{align*}
  which are linearly independent (entries $4$ and $5$ form $\mathrm{diag}(k,-k)$ with $k \ge 1$). For $\mathbf{x}^{(1)}$, the weighted sum over blocks $1,\dots,4$ is $k(-1) - 1 + 1 + k = 0$, so $(Q_{D^Q}\mathbf{x}^{(1)})_i = 0$ for all $i \ge 5$, while rows $1,2,3,4$ (with $\Tr_3 = 5k+6$, $\Tr_4 = 4k+5$) give:
  \begin{align*}
    (Q_{D^Q}\mathbf{x}^{(1)})_1 &= (5k+5)(-1) - 1 + 2 + k = -(4k+4), \\
    (Q_{D^Q}\mathbf{x}^{(1)})_2 &= -k - (4k+6) + 2 + k = -(4k+4), \\
    (Q_{D^Q}\mathbf{x}^{(1)})_3 &= -2k - 2 + (5k+6) + k = 4k+4, \\
    (Q_{D^Q}\mathbf{x}^{(1)})_4 &= -k - 1 + 1 + (4k+5)k = (4k+4)k.
  \end{align*}
  For $\mathbf{x}^{(2)}$, each of the $k+1$ tail pairs $(2r-1, 2r)$ ($3 \le r \le k+3$) contributes $2(-k)+k = -k$ to every earlier row $i < 2r-1$ and $c_i(-k+k) = 0$ to every later row $i > 2r$, while the weighted sum over blocks $1,\dots,4$ is $k(k+1)$. Checking rows $1,2,3,4$ and tail rows $2j-1, 2j$ ($3 \le j \le k+3$):
  \begin{align*}
    (Q_{D^Q}\mathbf{x}^{(2)})_1 &= (5k+5)(k+1) + (k+1) - 2(k+1) - k(k+1) = (4k+4)(k+1), \\
    (Q_{D^Q}\mathbf{x}^{(2)})_2 &= k(k+1) + (4k+6)(k+1) - 2(k+1) - k(k+1) = (4k+4)(k+1), \\
    (Q_{D^Q}\mathbf{x}^{(2)})_3 &= 2k(k+1) + 2(k+1) - (5k+6)(k+1) - k(k+1) = -(4k+4)(k+1), \\
    (Q_{D^Q}\mathbf{x}^{(2)})_4 &= k(k+1) + (k+1) - (k+1) - k(k+1) = 0, \\
    (Q_{D^Q}\mathbf{x}^{(2)})_{2j-1} &= 2k(k+1) + (5k+4+j)(-k) + k + (k+3-j)(-k) = (4k+4)(-k), \\
    (Q_{D^Q}\mathbf{x}^{(2)})_{2j} &= k(k+1) - k + (4k+7-j)k + (k+3-j)(-k) = (4k+4)k.
  \end{align*}
  Lifting $\mathbf{x}^{(1)}, \mathbf{x}^{(2)}$ to constant-on-blocks vectors in $\mathbb{R}^{3k+5}$ yields $\mult_{D^Q(G_k)}(4k+4) \ge 2$, with equality by Theorem~\ref{thm:CgraphDQ}.
\end{proof}

\begin{proposition}[Unique Minimality of $H_8 = C(1^8)$]\label{prop:minimality}
  Among all connected $\mathcal{C}$-graphs of order $n \le 8$, the $8$-vertex antiregular graph $G_1 = H_8 = C(1^8)$ is the unique counterexample to Conjecture~\ref{conj:MM} up to isomorphism.
\end{proposition}

\begin{proof}
  By~\cite{MM2025}, every connected $\mathcal{C}$-graph of order $n \ge 2$ is realized by an even-length positive integer composition $(\alpha_1,\dots,\alpha_m)$ of $n$ (and conversely every such composition produces a connected $\mathcal{C}$-graph). Compositions of $n$ of length $m$ correspond bijectively to $(m-1)$-element subsets of $\{1,\dots,n-1\}$; since a non-empty $(n-1)$-element set has exactly $2^{n-2}$ subsets of odd cardinality, there are precisely $2^{n-2}$ even-length creation sequences of order $n$, giving $\sum_{n=2}^{8} 2^{n-2} = 127$ creation sequences for $2 \le n \le 8$ (covering all connected $\mathcal{C}$-graphs of order $n \le 8$).

  For any connected $\mathcal{C}$-graph $G$, by~\cite[Theorem~4.3]{MM2026} the characteristic polynomial of $D^Q(G)$ splits into the local linear factors (with roots in $\mathcal{T}^Q(G)$) and $p(x) := \det(xI_m - Q_{D^Q}(G)) \in \mathbb{Z}[x]$. Hence, for every $\theta \notin \mathcal{T}^Q(G)$, $\mult_{D^Q(G)}(\theta)$ equals the algebraic multiplicity of $\theta$ as a root of $p(x)$, which equals the exponent $e$ of the unique irreducible factor $f(x) \in \mathbb{Q}[x]$ of $p(x)$ vanishing at $\theta$ (since irreducible polynomials over $\mathbb{Q}$ are separable). Factoring $p(x) = \prod_r f_r(x)^{e_r}$ into irreducible factors over $\mathbb{Q}[x]$ in SymPy across all $127$ creation sequences (zero floating-point arithmetic, zero root approximation) confirms that whenever $e_r > \nabla(G)$, either $f_r(x) = x - \theta$ with $\theta \in \mathcal{T}^Q(G)$, or $n = 8$ and the creation sequence is $(1^8)$. Since $(1^8)$ is the unique violating creation sequence, $H_8 = C(1^8)$ is the unique smallest counterexample up to isomorphism; its exact characteristic polynomial is
  \[
    \chi_{D^Q(H_8)}(x) = (x - 9)(x - 8)^2\bigl(x^5 - 55x^4 + 1139x^3 - 11261x^2 + 53580x - 98676\bigr),
  \]
  giving $\mult_{D^Q(H_8)}(8) = 2 > 1 = \nabla(H_8)$ with $\mathcal{T}^Q(H_8) = \emptyset$.
\end{proof}

\section{Universal Signless Laplacian Bound and Sharp Repaired Theorem for All $\mathcal{C}$-Graphs}\label{sec:Cgraph_bound}

\subsection{Block Recurrence, Critical-Value Separation, and Formal Tail Propagation}

Let $G = C(\alpha_1,\dots,\alpha_m)$ with $m = 2t$ even ($t \ge 1$). For $\mathbf{x} = (x_1,\dots,x_m)^T \in \mathbb{R}^m$, set
\[
  P_0 = 0, \qquad P_i = \sum_{j=1}^{i} \alpha_j x_j \quad (1 \le i \le m), \qquad
  S_{2r} = \sum_{j=r}^{t} \alpha_{2j} x_{2j} \quad (1 \le r \le t), \qquad S_{2t+2} = 0,
\]
and define the \emph{critical values}
\begin{equation}\label{eq:crit}
  c_{2r-1} := d_{2r-1} - 1 + \alpha_{2r-1}, \qquad
  c_{2r} := d_{2r} - \alpha_{2r} \qquad (1 \le r \le t).
\end{equation}

\begin{lemma}[Exact Block Recurrence and Conservation Law]\label{lem:recurrence}
  A vector $\mathbf{x} \in \mathbb{R}^m$ satisfies $Q_\pi(G)\mathbf{x} = \mu \mathbf{x}$ if and only if, for all $r = 1,\dots,t$,
  \begin{align}
    (\mu - c_{2r-1}) x_{2r-1} &= S_{2r}, \label{eq:rec_odd}\\
    (\mu - c_{2r}) x_{2r} &= P_{2r-1} + S_{2r}, \label{eq:rec_even}
  \end{align}
  with $P_{2r-1} = P_{2r-2} + \alpha_{2r-1} x_{2r-1}$, $P_{2r} = P_{2r-1} + \alpha_{2r} x_{2r}$, $S_{2r+2} = S_{2r} - \alpha_{2r} x_{2r}$, and $S_{2t+2} = 0$. Moreover, for every $r = 1,\dots,t$,
  \begin{equation}\label{eq:conservation}
    P_{2r} + S_{2r+2} = P_{2r-1} + S_{2r}.
  \end{equation}
\end{lemma}

\begin{proof}
  Immediate from~\eqref{eq:Qpi}: $(Q_\pi \mathbf{x})_{2r-1} = c_{2r-1} x_{2r-1} + S_{2r}$ and $(Q_\pi \mathbf{x})_{2r} = d_{2r} x_{2r} + P_{2r-1} + (S_{2r} - \alpha_{2r} x_{2r}) = c_{2r} x_{2r} + P_{2r-1} + S_{2r}$. Adding $P_{2r} = P_{2r-1} + \alpha_{2r} x_{2r}$ and $S_{2r+2} = S_{2r} - \alpha_{2r} x_{2r}$ yields~\eqref{eq:conservation}.
\end{proof}

\begin{lemma}[Strict Separation of Critical Values]\label{lem:separation}
  Let $O_1 = \{j \in \{3,\dots,m\} : \alpha_j = 1\}$. Then:
  \begin{enumerate}
    \item $c_1 - c_2 = \alpha_1 + 2\alpha_2 - 2 \ge 1$, so $c_1 > c_2$.
    \item On $O_1$, $c_{2r-1}$ is strictly decreasing in $r$, $c_{2s}$ is strictly increasing in $s$, and $c_{2s} - c_{2r-1} \ge 1$ for every odd $2r-1 \in O_1$ and even $2s \in O_1$. Thus $\{c_j : j \in O_1\}$ are pairwise distinct.
    \item If $\alpha_2 = 1$, then $c_{2s} > c_2$ for every even $2s \in O_1$, and $c_{2r-1} \le c_2$ for every odd $2r-1 \in O_1$, with equality only when $\alpha_1 = 1$ and $2r-1 = 3$.
  \end{enumerate}
\end{lemma}

\begin{proof}
  Using $d_1 = (\alpha_1-1) + \sum_{j=1}^t \alpha_{2j}$, $d_2 = \alpha_1 + \sum_{j=2}^t \alpha_{2j}$, $d_{2r-1} = (\alpha_{2r-1}-1) + \sum_{j=r}^t \alpha_{2j}$, and $d_{2s} = \sum_{i=1}^{2s-1}\alpha_i + \sum_{j=s+1}^t \alpha_{2j}$:
  \begin{enumerate}
    \item $c_1 - c_2 = (2\alpha_1 - 2 + \alpha_2) - (\alpha_1 - \alpha_2) = \alpha_1 + 2\alpha_2 - 2 \ge 1$.
    \item For odd $2r-1 \in O_1$, $c_{2r-1} = \sum_{j=r}^t \alpha_{2j}$ strictly decreases in $r$. For even $2s \in O_1$, $c_{2s} = \sum_{i=1}^{2s-1}\alpha_i + \sum_{j=s+1}^t \alpha_{2j} - 1$, so for $2s_1 < 2s_2$ in $O_1$, $c_{2s_2} - c_{2s_1} = \alpha_{2s_1} + \sum_{k=s_1}^{s_2-1}\alpha_{2k+1} - \alpha_{2s_2} = \sum_{k=s_1}^{s_2-1}\alpha_{2k+1} \ge 1$ (using $\alpha_{2s_1} = \alpha_{2s_2} = 1$), strictly increasing in $s$. Comparing odd $2r-1 \in O_1$ and even $2s \in O_1$ ($r,s \ge 2$): if $s < r$, then $c_{2s} - c_{2r-1} = \sum_{i=1}^{2s-1}\alpha_i + \sum_{j=s+1}^{r-1}\alpha_{2j} - 1 \ge 2$; if $s \ge r$, cancelling $\{2j : r \le j \le s-1\} \subset \{1,\dots,2s-1\}$ leaves $c_{2s} - c_{2r-1} \ge (\alpha_1 + \alpha_2 + \alpha_{2s-1}) - \alpha_{2s} - 1 \ge 1$.
    \item If $\alpha_2 = 1$, then $c_2 = (\alpha_1 - 1) + \sum_{j=2}^t \alpha_{2j}$, so $c_2 - c_{2r-1} = (\alpha_1 - 1) + \sum_{j=2}^{r-1}\alpha_{2j} \ge 0$ (vanishing iff $\alpha_1=1, r=2$) and $c_{2s} - c_2 \ge (\alpha_2 + \alpha_{2s-1}) - \alpha_{2s} = \alpha_{2s-1} \ge 1$. \qedhere
  \end{enumerate}
\end{proof}

To avoid conflating formal recurrence variables with actual eigenvectors, we formulate tail propagation strictly for \emph{formal recurrence sequences} without imposing the terminal boundary condition $\tilde{S}_{2t+2} = 0$.

\begin{lemma}[Formal Recurrence Tail Propagation]\label{lem:sign}
  Let $\mu \in \mathbb{R}$ and $i^* \in \{2,\dots,m\}$ ($m = 2t \ge 2$) satisfy that for every index $j$ with $i^* < j \le m$, $\mu - c_j > 0$ when $j$ is odd and $\mu - c_j < 0$ when $j$ is even. Consider a formal recurrence sequence $(\tilde{x}_j, \tilde{P}_j, \tilde{S}_{2\ell})$ generated for $j > i^*$ by the update rules
  \begin{align*}
    \tilde{x}_{2\ell-1} &= \frac{\tilde{S}_{2\ell}}{\mu - c_{2\ell-1}}, & \tilde{P}_{2\ell-1} &= \tilde{P}_{2\ell-2} + \alpha_{2\ell-1}\tilde{x}_{2\ell-1}, \\
    \tilde{x}_{2\ell} &= \frac{\tilde{P}_{2\ell-1} + \tilde{S}_{2\ell}}{\mu - c_{2\ell}}, & \tilde{P}_{2\ell} &= \tilde{P}_{2\ell-1} + \alpha_{2\ell}\tilde{x}_{2\ell}, \qquad \tilde{S}_{2\ell+2} = \tilde{S}_{2\ell} - \alpha_{2\ell}\tilde{x}_{2\ell},
  \end{align*}
  without imposing $\tilde{S}_{2t+2} = 0$. Suppose one of the following initial conditions holds:
  \begin{enumerate}
    \item[(a)] $i^* = 2r - 1$ ($2 \le r \le t$), $\tilde{P}_{2r-2} = 0$, $\tilde{S}_{2r} = 0$, and $\tilde{x}_{2r-1} = 1$ (so $\tilde{P}_{2r-1} = \alpha_{2r-1}$);
    \item[(b)] $i^* = 2r$ ($2 \le r \le t$), $\tilde{P}_{2r-1} = 0$, $\tilde{S}_{2r} = 0$, and $\tilde{x}_{2r} = -1$ (so $\tilde{P}_{2r} = -\alpha_{2r}$, $\tilde{S}_{2r+2} = \alpha_{2r}$);
    \item[(c)] $i^* = 2$, $\tilde{P}_2 = -1$, and $\tilde{S}_4 = 1$.
  \end{enumerate}
  Then the terminal formal value satisfies $\tilde{S}_{2t+2} > 0$.
\end{lemma}

\begin{proof}
  Define $\tilde{T}_\ell := \tilde{P}_{2\ell-1} + \tilde{S}_{2\ell}$ for $1 \le \ell \le t$. Note that the formal update rules preserve the conservation identity $\tilde{P}_{2\ell} + \tilde{S}_{2\ell+2} = \tilde{P}_{2\ell-1} + \tilde{S}_{2\ell} = \tilde{T}_\ell$ for every $\ell$.

  \emph{Inductive step.} Suppose for some $\ell \in \{1,\dots,t\}$ with $2\ell > i^*$ we have $\tilde{S}_{2\ell} \ge 0$ and $\tilde{T}_\ell > 0$. Since $\mu - c_{2\ell} < 0$, we obtain $\tilde{x}_{2\ell} = \tilde{T}_\ell / (\mu - c_{2\ell}) < 0$, whence
  \[
    \tilde{S}_{2\ell+2} = \tilde{S}_{2\ell} - \alpha_{2\ell}\tilde{x}_{2\ell} > \tilde{S}_{2\ell} \ge 0, \qquad
    \tilde{P}_{2\ell} + \tilde{S}_{2\ell+2} = \tilde{T}_\ell > 0.
  \]
  If $\ell = t$, we conclude $\tilde{S}_{2t+2} > 0$. If $\ell < t$, then $2\ell + 1 > i^*$ is odd so $\mu - c_{2\ell+1} > 0$, giving $\tilde{x}_{2\ell+1} = \tilde{S}_{2\ell+2} / (\mu - c_{2\ell+1}) > 0$ and
  \[
    \tilde{T}_{\ell+1} = \tilde{P}_{2\ell+1} + \tilde{S}_{2\ell+2} = (\tilde{P}_{2\ell} + \tilde{S}_{2\ell+2}) + \alpha_{2\ell+1}\tilde{x}_{2\ell+1} = \tilde{T}_\ell + \alpha_{2\ell+1}\tilde{x}_{2\ell+1} > \tilde{T}_\ell > 0.
  \]
  By induction on $\ell$, we reach $\tilde{S}_{2t+2} > 0$.

  \emph{Verification of initial conditions (including boundary cases $i^* = m$ and $t = 1$).}
  \begin{itemize}
    \item In~(a) ($i^* = 2r - 1 < 2r$): $\tilde{S}_{2r} = 0 \ge 0$ and $\tilde{T}_r = \tilde{P}_{2r-1} + \tilde{S}_{2r} = \alpha_{2r-1} > 0$, entering the induction at $\ell = r$.
    \item In~(b) ($i^* = 2r$): we have $\tilde{S}_{2r+2} = \alpha_{2r} > 0$ and $\tilde{P}_{2r} + \tilde{S}_{2r+2} = 0$. If $i^* = m$ (i.e., $r = t$), then $\tilde{S}_{2t+2} = \alpha_{2t} > 0$ holds immediately. If $r < t$, then $2r + 1 > i^*$ gives $\tilde{x}_{2r+1} = \tilde{S}_{2r+2}/(\mu - c_{2r+1}) > 0$ and $\tilde{T}_{r+1} = (\tilde{P}_{2r} + \tilde{S}_{2r+2}) + \alpha_{2r+1}\tilde{x}_{2r+1} = \alpha_{2r+1}\tilde{x}_{2r+1} > 0$, entering the induction at $\ell = r + 1$.
    \item In~(c) ($i^* = 2$): if $t = 1$ (so $m = 2$), then $\tilde{S}_{2t+2} = \tilde{S}_4 = 1 > 0$ holds immediately. If $t \ge 2$, then $3 > i^*$ gives $\tilde{x}_3 = \tilde{S}_4 / (\mu - c_3) > 0$ and $\tilde{T}_2 = (\tilde{P}_2 + \tilde{S}_4) + \alpha_3 \tilde{x}_3 = (-1 + 1) + \alpha_3 \tilde{x}_3 = \alpha_3 \tilde{x}_3 > 0$ with $\tilde{S}_4 = 1 > 0$, entering the induction at $\ell = 2$. \qedhere
  \end{itemize}
\end{proof}

\subsection{Main Multiplicity Theorems for All Connected $\mathcal{C}$-Graphs}

\begin{theorem}[Universal Signless Laplacian Dilworth Bound for $\mathcal{C}$-Graphs]\label{thm:CgraphQ}
  For every connected $\mathcal{C}$-graph $G = C(\alpha_1,\dots,\alpha_m)$ ($m = 2t \ge 2$ even) and every $\mu \in \mathbb{R}$, $\mult_{Q_\pi(G)}(\mu) \le \nabla(G)$. In particular, for every $\mu \notin \mathcal{T}_Q(G)$,
  \[
    \mult_{Q(G)}(\mu) \le \nabla(G).
  \]
\end{theorem}

\begin{proof}
  \emph{Base case $m = 2$ ($t = 1$).} Here $\nabla(G) = 1$ and~\eqref{eq:Qpi} gives
  \[
    Q_\pi(G) = \begin{pmatrix} 2\alpha_1 + \alpha_2 - 2 & \alpha_2 \\ \alpha_1 & \alpha_1 \end{pmatrix}.
  \]
  Its discriminant is $(\alpha_1 + \alpha_2 - 2)^2 + 4\alpha_1\alpha_2 > 0$, so its two eigenvalues are distinct and $\mult_{Q_\pi(G)}(\mu) \le 1 = \nabla(G)$.

  Assume henceforth that $m = 2t \ge 4$ ($t \ge 2$). Let $S = \{j \ge 3 : \alpha_j \ge 2\}$, $O_1 = \{j \ge 3 : \alpha_j = 1\}$, and $Z(\mu) = \{i \in \{1,\dots,m\} : c_i = \mu\}$, so $\nabla(G) = 1 + |S|$ by Theorem~\ref{thm:dilworth}, and $|Z(\mu) \cap \{1,2\}| \le 1$, $|Z(\mu) \cap O_1| \le 1$ by Lemma~\ref{lem:separation}. By Lemma~\ref{lem:diag}, $\mult_{Q_\pi(G)}(\mu) = \dim\ker(Q_\pi(G) - \mu I_m)$.

  \emph{Free-parameter isomorphism.} Define a linear map $\Phi : \mathbb{R}^{1 + |Z(\mu)|} \to \mathbb{R}^m$ taking a parameter tuple $p = \bigl(s_2, \{z_i : i \in Z(\mu)\}\bigr)$ to $\tilde{\mathbf{x}} = (\tilde{x}_1,\dots,\tilde{x}_m)^T$ via the formal recurrence initialized at $\tilde{P}_0 := 0$ and $\tilde{S}_2 := s_2$: for $r = 1,\dots,t$,
  \begin{itemize}
    \item if $2r-1 \notin Z(\mu)$, set $\tilde{x}_{2r-1} := \tilde{S}_{2r}/(\mu - c_{2r-1})$; if $2r-1 \in Z(\mu)$, set $\tilde{x}_{2r-1} := z_{2r-1}$ and define the linear functional $L_{2r-1}(p) := \tilde{S}_{2r}$. Then set $\tilde{P}_{2r-1} := \tilde{P}_{2r-2} + \alpha_{2r-1}\tilde{x}_{2r-1}$;
    \item if $2r \notin Z(\mu)$, set $\tilde{x}_{2r} := (\tilde{P}_{2r-1} + \tilde{S}_{2r})/(\mu - c_{2r})$; if $2r \in Z(\mu)$, set $\tilde{x}_{2r} := z_{2r}$ and define the linear functional $L_{2r}(p) := \tilde{P}_{2r-1} + \tilde{S}_{2r}$. Then set $\tilde{P}_{2r} := \tilde{P}_{2r-1} + \alpha_{2r}\tilde{x}_{2r}$ and $\tilde{S}_{2r+2} := \tilde{S}_{2r} - \alpha_{2r}\tilde{x}_{2r}$.
  \end{itemize}
  Finally define $L_{\mathrm{end}}(p) := \tilde{S}_{2t+2}$, and assemble the linear constraint map
  \[
    C : \mathbb{R}^{1 + |Z(\mu)|} \longrightarrow \mathbb{R}^{1 + |Z(\mu)|}, \qquad
    C(p) := \bigl(\{L_i(p)\}_{i \in Z(\mu)},\, L_{\mathrm{end}}(p)\bigr).
  \]
  We verify that $\Phi|_{\ker C} : \ker C \to \ker(Q_\pi - \mu I_m)$ is a linear isomorphism:
  \begin{itemize}
    \item If $p \in \ker C$, then $\tilde{S}_{2t+2} = L_{\mathrm{end}}(p) = 0$; backward induction on $\tilde{S}_{2r} = \tilde{S}_{2r+2} + \alpha_{2r}\tilde{x}_{2r}$ yields $\tilde{S}_{2r} = \sum_{j=r}^t \alpha_{2j}\tilde{x}_{2j}$ for all $r$. Every non-critical equation holds by definition of $\tilde{x}_j$, and every critical equation holds because $L_i(p) = 0$ for $i \in Z(\mu)$. By Lemma~\ref{lem:recurrence}, $\Phi(p) \in \ker(Q_\pi - \mu I_m)$.
    \item Conversely, given $\mathbf{x} \in \ker(Q_\pi - \mu I_m)$, set $p(\mathbf{x}) := \bigl(\sum_{j=1}^t \alpha_{2j} x_{2j},\, \{x_i\}_{i \in Z(\mu)}\bigr)$. By Lemma~\ref{lem:recurrence}, induction on $r$ shows that $\Phi(p(\mathbf{x})) = \mathbf{x}$, $\tilde{S}_{2r} = S_{2r}$, and $C(p(\mathbf{x})) = \mathbf{0}$. Moreover, $\Phi(p) = \mathbf{0}$ forces $z_i = \tilde{x}_i = 0$ ($i \in Z(\mu)$) and $s_2 = \tilde{S}_2 = \sum_{j=1}^t \alpha_{2j}\tilde{x}_{2j} = 0$ when $p \in \ker C$, so $\Phi|_{\ker C}$ is injective.
  \end{itemize}
  By the rank-nullity theorem,
  \begin{equation}\label{eq:dim_formula}
    \mult_{Q_\pi(G)}(\mu) = \dim\ker(Q_\pi - \mu I_m) = \dim\ker C = 1 + |Z(\mu)| - \rank(C).
  \end{equation}

  \emph{Case 1: $\alpha_2 = 1$ and $2 \in Z(\mu)$ ($\mu = c_2$, so $1 \notin Z(\mu)$ by Lemma~\ref{lem:separation}(1)).}
  Since $\mu - c_1 = c_2 - c_1 = -\alpha_1$, we have $\tilde{x}_1 = -s_2/\alpha_1$ and $\tilde{P}_1 = -s_2$, so $L_2(p) = \tilde{P}_1 + \tilde{S}_2 = 0$ vanishes identically. Changing variables to $(s_2, y, \{z_i : i \in Z'(\mu)\})$ where $y := s_2 - z_2$ and $Z'(\mu) := Z(\mu) \cap \{3,\dots,m\}$, we obtain $\tilde{P}_2 = -y$ and $\tilde{S}_4 = y$. Thus $s_2$ does not appear in any tail equation ($j \ge 3$) or in $L_{\mathrm{end}}$, spanning only the $1$-dimensional kernel direction $(-1/\alpha_1, 1, 0,\dots,0)^T$, while the remaining tail functionals $C_{\mathrm{tail}} : \mathbb{R}^{1 + |Z'(\mu)|} \to \mathbb{R}^{1 + |Z'(\mu)|}$ act on $p_{\mathrm{tail}} = (y, \{z_i : i \in Z'(\mu)\})$, giving
  \[
    \mult_{Q_\pi(G)}(\mu) = 1 + \bigl(1 + |Z'(\mu)| - \rank(C_{\mathrm{tail}})\bigr).
  \]
  By Lemma~\ref{lem:separation}(3), $Z'(\mu) \cap O_1 \subseteq \{3\}$, so $|Z'(\mu)| \le |S| + 1$:
  \begin{itemize}
    \item If $|Z'(\mu)| \le |S| - 1$, then $1 + (1 + |Z'(\mu)| - \rank(C_{\mathrm{tail}})) \le 1 + |S| = \nabla(G)$ directly.
    \item If $|Z'(\mu)| = |S| + 1$, then $S \subseteq Z'(\mu)$ and $3 \in Z'(\mu) \cap O_1$. The row-$3$ constraint is $L_3(p_{\mathrm{tail}}) = \tilde{S}_4 = y$, which is a non-zero coordinate functional ($L_3(1, \mathbf{0}) = 1$). Let $i_{\max} := \max Z'(\mu) \ge 3$. Since $S \subseteq Z'(\mu)$, every $j > i_{\max}$ lies in $O_1 \setminus \{3\}$, so $\mu - c_j > 0$ for odd $j > i_{\max}$ and $\mu - c_j < 0$ for even $j > i_{\max}$ by Lemma~\ref{lem:separation}(3). Define the test vector $p^* \in \mathbb{R}^{1+|Z'(\mu)|}$ by $y = 0$, $z_i = 0$ ($i \in Z'(\mu) \setminus \{i_{\max}\}$), and $z_{i_{\max}} = 1$ (if $i_{\max}$ is odd) or $-1$ (if $i_{\max}$ is even). Then $L_3(p^*) = 0$, while Lemma~\ref{lem:sign}(a)--(b) gives $L_{\mathrm{end}}(p^*) = \tilde{S}_{2t+2} > 0$. Hence $L_3$ and $L_{\mathrm{end}}$ are linearly independent, so $\rank(C_{\mathrm{tail}}) \ge 2$ and $\mult_{Q_\pi}(\mu) \le 1 + (1 + |S| + 1 - 2) = 1 + |S| = \nabla(G)$.
    \item If $|Z'(\mu)| = |S|$: if $3 \in Z'(\mu)$, then $L_3(p_{\mathrm{tail}}) = y \not\equiv 0$ gives $\rank(C_{\mathrm{tail}}) \ge 1$. Otherwise $3 \notin Z'(\mu)$, so $Z'(\mu) = S$. If $S \ne \emptyset$, the test vector $p^*$ at $i_{\max} = \max S \ge 3$ gives $L_{\mathrm{end}}(p^*) > 0$ by Lemma~\ref{lem:sign}(a)--(b). If $S = \emptyset$, then $Z'(\mu) = \emptyset$ and every $j \in \{3,\dots,m\}$ lies in $O_1$; because $m \ge 4$ and $3 \notin Z'(\mu)$, we have $c_3 < c_2 = \mu$ (so $\alpha_1 \ge 2$ by Lemma~\ref{lem:separation}(3)), hence $\mu - c_j > 0$ for odd $j \ge 3$ and $\mu - c_j < 0$ for even $j \ge 4$. Setting $y = 1$ gives $\tilde{P}_2 = -1, \tilde{S}_4 = 1$, and Lemma~\ref{lem:sign}(c) yields $L_{\mathrm{end}}(1) = \tilde{S}_{2t+2} > 0$. In all subcases $\rank(C_{\mathrm{tail}}) \ge 1$, giving $\mult_{Q_\pi}(\mu) \le 1 + (1 + |S| - 1) = \nabla(G)$.
  \end{itemize}

  \emph{Case 2: $2 \notin Z(\mu)$ or $\alpha_2 \ge 2$.}
  Whenever $i_0 \in Z(\mu) \cap \{1,2\}$, the functional $L_{i_0}(p)$ is a non-zero scalar multiple of $s_2$: directly $L_1(p) = s_2$ if $i_0 = 1$; and if $i_0 = 2$ with $\alpha_2 \ge 2$, then $\mu - c_1 = c_2 - c_1 = -(\alpha_1 + 2\alpha_2 - 2) < 0$ gives
  \[
    L_2(p) = \tilde{P}_1 + \tilde{S}_2 = \left(\frac{\alpha_1}{c_2 - c_1} + 1\right) s_2 = \frac{2\alpha_2 - 2}{\alpha_1 + 2\alpha_2 - 2}\,s_2,
  \]
  whose coefficient is strictly positive. Since $|Z(\mu)| \le |S| + 2$:
  \begin{itemize}
    \item If $|Z(\mu)| \le |S|$, then~\eqref{eq:dim_formula} gives $\mult_{Q_\pi}(\mu) \le 1 + |S| = \nabla(G)$.
    \item If $|Z(\mu)| = |S| + 1$ and $Z(\mu) \cap \{1,2\} = \{i_0\}$, then $L_{i_0} \not\equiv 0$ gives $\rank(C) \ge 1$, so $\mult_{Q_\pi}(\mu) \le 1 + (|S|+1) - 1 = \nabla(G)$.
    \item If $|Z(\mu)| = |S| + 1$ and $Z(\mu) \cap \{1,2\} = \emptyset$, then $S \subseteq Z(\mu)$ and $Z(\mu) \cap O_1 = \{j_0\}$. Set $i_{\max} := \max Z(\mu) \ge 3$. Every $j > i_{\max}$ lies in $O_1$, and since $j_0 \le i_{\max}$, Lemma~\ref{lem:separation}(2) guarantees $\mu - c_j > 0$ for odd $j > i_{\max}$ and $\mu - c_j < 0$ for even $j > i_{\max}$. The test vector $p^*$ with $s_2 = 0$, $z_i = 0$ ($i \ne i_{\max}$), and $z_{i_{\max}} = \pm 1$ yields $L_{\mathrm{end}}(p^*) > 0$ by Lemma~\ref{lem:sign}(a)--(b), so $\rank(C) \ge 1$.
    \item If $|Z(\mu)| = |S| + 2$, then $Z(\mu) \cap \{1,2\} = \{i_0\}$, $S \subseteq Z(\mu)$, and $Z(\mu) \cap O_1 = \{j_0\}$. Set $i_{\max} := \max Z(\mu) \ge 3$. The functional $L_{i_0}$ is a non-zero multiple of $s_2$, so $L_{i_0}(1, \mathbf{0}) \ne 0$; meanwhile, the test vector $p^*$ with $s_2 = 0$, $z_i = 0$ ($i \in Z(\mu) \setminus \{i_{\max}\}$), and $z_{i_{\max}} = \pm 1$ satisfies $L_{i_0}(p^*) = 0$ and (since every $j > i_{\max}$ lies in $O_1$ and satisfies the sign condition by Lemma~\ref{lem:separation}(2), as in the preceding subcase) $L_{\mathrm{end}}(p^*) > 0$ by Lemma~\ref{lem:sign}(a)--(b). Hence $L_{i_0}$ and $L_{\mathrm{end}}$ are linearly independent, so $\rank(C) \ge 2$ and $\mult_{Q_\pi}(\mu) \le 1 + (|S|+2) - 2 = \nabla(G)$.
  \end{itemize}
  Finally, if $\mu \notin \mathcal{T}_Q(G)$, then $\mult_{Q(G)}(\mu) = \mult_{Q_\pi(G)}(\mu) \le \nabla(G)$.
\end{proof}

\begin{theorem}[Sharp Repaired Bound for Mandal--Manna Conjecture~\ref{conj:MM}]\label{thm:CgraphDQ}
  For every connected $\mathcal{C}$-graph $G$ and every $\theta \notin \mathcal{T}^Q(G)$,
  \[
    \mult_{D^Q(G)}(\theta) \le \nabla(G) + 1.
  \]
  This bound is sharp, attained by $G_k$ ($k \ge 1$) of Theorem~\ref{thm:disproof}.
\end{theorem}

\begin{proof}
  Let $\theta \notin \mathcal{T}^Q(G)$. By Lemma~\ref{lem:rank1}(3), $\mu := 2n - 4 - \theta \notin \mathcal{T}_Q(G)$, so Theorem~\ref{thm:CgraphQ} gives $\mult_{Q(G)}(\mu) \le \nabla(G)$. Applying Lemma~\ref{lem:rank1}(2) yields $\mult_{D^Q(G)}(\theta) \le \mult_{Q(G)}(\mu) + 1 \le \nabla(G) + 1$.
\end{proof}

\section{General Cographs: Universal Linear Upper Bounds and Linear Separation from $\nabla(G)$}\label{sec:cographs}

For any cograph $G$ with integer vertex weight $w : V(G) \to \mathbb{Z}$, write $Q_w(G) := Q(G) + \mathrm{diag}(w)$. Call $u, v \in V(G)$ \emph{$(G,w)$-twins} if they are twins in $G$ and $w(u) = w(v)$; let $\mathcal{T}_{Q,w}(G)$ denote the set of local eigenvalues of $Q_w(G)$ arising from maximal $(G,w)$-twin classes $C$ of size $|C| \ge 2$ (namely $d_C - 1 + w_C$ if $C$ is a clique, and $d_C + w_C$ if $C$ is an independent set). We say $w$ is \emph{vicinal-monotone} on $G$ if $u \preceq_G v \implies w(u) \le w(v)$ for all $u, v \in V(G)$ (note that if $u, v$ are twins in $G$, then $u \preceq_G v$ and $v \preceq_G u$, so any vicinal-monotone weight on $G$ is automatically constant on every maximal twin class of $G$).

\subsection{Cotree Dilworth Formula and Weighted Threshold Simplicity}

\begin{lemma}[Unified Cotree Formula for $\nabla(G)$]\label{lem:cotree_nabla}
  Let $G$ be a finite cograph.
  \begin{enumerate}
    \item If $G$ is complete or edgeless (in particular, if $G = K_1$), then $\nabla(G) = 1$.
    \item If $G = G_1 \cup \dots \cup G_r \cup \overline{K_a}$ is disconnected, where $G_1,\dots,G_r$ ($r \ge 0$) are its nontrivial connected components ($|G_i| \ge 2$) and $a \ge 0$ is the number of isolated vertices, then
    \[
      \nabla(G) = \max\Bigl(1,\, \sum_{i=1}^{r} \nabla(G_i)\Bigr).
    \]
    \item If $G = K_b \vee H_1 \vee \dots \vee H_r$ is connected, where $H_1,\dots,H_r$ ($r \ge 0$) are its nontrivial disconnected co-connected parts ($|H_i| \ge 2$) and $b \ge 0$ is the number of universal vertices, then
    \[
      \nabla(G) = \max\Bigl(1,\, \sum_{i=1}^{r} \nabla(H_i)\Bigr).
    \]
  \end{enumerate}
\end{lemma}

\begin{proof}
  Part~(1) is immediate since all vertices are mutually twins.

  For~(2), if $r = 0$ then $G = \overline{K_a}$ is edgeless and $\nabla(G) = 1 = \max(1, 0)$. Assume $r \ge 1$. Every isolated vertex $v$ satisfies $N(v) = \emptyset \subseteq N[u]$ for all $u \in V(G)$ and can be placed at the bottom of any chain, so deleting $\overline{K_a}$ does not change $\nabla(G)$. For $1 \le i < j \le r$, any $u \in G_i$ and $v \in G_j$ are non-adjacent; because $G_i$ and $G_j$ are connected on $\ge 2$ vertices, $u$ has a neighbor in $G_i \setminus N[v]$ and $v$ has a neighbor in $G_j \setminus N[u]$, so $u$ and $v$ are incomparable. Within each $G_i$, the vicinal preorder agrees with that of $G$. Hence maximum antichains partition across $G_1,\dots,G_r$, giving $\nabla(G) = \sum_{i=1}^r \nabla(G_i) = \max(1, \sum_{i=1}^r \nabla(G_i))$.

  For~(3), if $r = 0$ then $G = K_b$ is complete and $\nabla(G) = 1 = \max(1, 0)$. Assume $r \ge 1$. Every universal vertex $v \in K_b$ satisfies $N(u) \subseteq V(G) = N[v]$ for all $u$ and sits at the top of any chain, so deleting $K_b$ does not change $\nabla(G)$. For $1 \le i < j \le r$, any $u \in H_i$ and $v \in H_j$ are adjacent; because $H_i$ is disconnected, $u$ is non-adjacent to some $u' \in H_i$ whereas $v$ is adjacent to $u'$, so $v \not\preceq u$, and symmetrically $u \not\preceq v$. Within each $H_i$, every pair $u, u' \in H_i$ shares identical adjacencies outside $H_i$, so $N_G(u) \subseteq N_G[u'] \iff N_{H_i}(u) \subseteq N_{H_i}[u']$. Hence $\nabla(G) = \sum_{i=1}^r \nabla(H_i) = \max(1, \sum_{i=1}^r \nabla(H_i))$.
\end{proof}

\begin{lemma}[Canonical Block Form of Weighted Connected Threshold Graphs]\label{lem:thresh_canon}
  Let $G$ be a connected threshold graph on $n \ge 2$ vertices, and let $w : V(G) \to \mathbb{Z}$ be an integer weight function that is non-decreasing in the vicinal preorder of $G$ (or, more generally, constant on each block of an alternating threshold decomposition and non-decreasing along the chain $A_s \preceq \dots \preceq A_1 \preceq B_1 \preceq \dots \preceq B_s$). Then $V(G)$ admits an alternating partition into independent blocks $A_1,\dots,A_s$ ($|A_i|=a_i$) and clique blocks $B_1,\dots,B_s$ ($|B_i|=b_i$, $s \ge 1$), where $A_i \sim B_j \iff i \le j$, all $B_i$ are pairwise adjacent, all $A_i$ are pairwise non-adjacent, and $w$ is constant on each block with $w_{A_s} \le \dots \le w_{A_{i_0}} \le w_{B_1} \le \dots \le w_{B_s}$, falling into one of four mutually exclusive cases:
  \begin{enumerate}
    \item[(I)] $a_1 \ge 2$, and all other $a_i, b_i \ge 1$ ($i_0 = 1$);
    \item[(II)] $a_1 = 0$ ($A_1$ absent), $b_1 \ge 2$, and all other $a_i, b_i \ge 1$ ($i_0 = 2$);
    \item[(III)] $a_1 = 1$, all other $a_i, b_i \ge 1$, and $w_{B_1} - w_{A_1} \ge 1$ ($i_0 = 1$);
    \item[(IV)] $a_1 = 0$ ($A_1$ absent), $b_1 = 1$, $s \ge 2$, all other $a_i, b_i \ge 1$, and $w_{B_1} - w_{A_2} \ge 1$ ($i_0 = 2$).
  \end{enumerate}
\end{lemma}

\begin{proof}
  Every connected threshold graph on $n \ge 2$ vertices is built by iteratively adding universal vertices and isolated vertices, terminating with a non-empty universal clique $B_s$, which yields an alternating partition $A_1, B_1, \dots, A_s, B_s$ ($s \ge 1$) with $a_2,\dots,a_s \ge 1$, $b_1,\dots,b_s \ge 1$, and $a_1 \ge 0$ (where $a_1 = 0$ means $A_1$ is absent).
  \begin{itemize}
    \item If $a_1 \ge 2$, we are in case~(I).
    \item If $a_1 = 0$ and $b_1 \ge 2$, we are in case~(II).
    \item If $a_1 = 0$ and $b_1 = 1$, then $s \ge 2$ (as $n \ge 2$), and the single vertex $v \in B_1$ has open neighborhood $N_G(v) = B_2 \cup \dots \cup B_s = N_G(u)$ for every $u \in A_2$, so $B_1 \cup A_2$ is an independent twin class in $G$. By monotonicity, $w_{A_2} \le w_{B_1}$: if $w_{A_2} = w_{B_1}$, then $B_1$ and $A_2$ are $(G,w)$-twins and merging $B_1$ into $A_2$ re-indexes $s \mapsto s-1$ with $a_1' = a_2 + 1 \ge 2$, placing $(G,w)$ in case~(I); if $w_{B_1} > w_{A_2}$, then since $w$ is integer-valued, $w_{B_1} - w_{A_2} \ge 1$, placing $(G,w)$ in case~(IV).
    \item If $a_1 = 1$, then the single vertex $u \in A_1$ has closed neighborhood $N_G[u] = \{u\} \cup B_1 \cup \dots \cup B_s = N_G[v]$ for every $v \in B_1$, so $A_1 \cup B_1$ is a clique twin class in $G$. Since $w_{A_1} \le w_{B_1}$ and $w$ is integer-valued: if $w_{A_1} = w_{B_1}$, then $A_1$ and $B_1$ are $(G,w)$-twins and merging $A_1$ into $B_1$ gives $a_1' = 0$ and $b_1' = b_1 + 1 \ge 2$, placing $(G,w)$ in case~(II); if $w_{B_1} > w_{A_1}$, then $w_{B_1} - w_{A_1} \ge 1$, placing $(G,w)$ in case~(III). \qedhere
  \end{itemize}
\end{proof}

\begin{theorem}[Weighted Signless Laplacian Simplicity for Connected Threshold Graphs]\label{thm:Qsimple}
  Let $G$ be a connected threshold graph on $n \ge 1$ vertices with an integer weight $w : V(G) \to \mathbb{Z}$ satisfying the hypotheses of Lemma~\ref{lem:thresh_canon} (or $n = 1$). Then every eigenvalue $\mu \notin \mathcal{T}_{Q,w}(G)$ of $Q_w(G)$ is simple:
  \[
    \mult_{Q_w(G)}(\mu) \le 1.
  \]
\end{theorem}

\begin{proof}
  If $n = 1$ ($G = K_1$), then $\dim \mathbb{R}^1 = 1$ and the claim is trivial. Assume $n \ge 2$ and place $(G, w)$ in its canonical block form~(I), (II), (III), or~(IV) from Lemma~\ref{lem:thresh_canon}.

  \emph{Complete-graph boundary cases ($s = 1$ in cases~(II) and~(III)).}
  \begin{itemize}
    \item In case~(II) with $s = 1$, $G = K_n$ ($n = b_1 \ge 2$) with constant weight $w_{B_1}$, so $\spec(Q_w(K_n)) = \{2n-2+w_{B_1},\, (n-2+w_{B_1})^{(n-1)}\}$ and $\mathcal{T}_{Q,w}(K_n) = \{n-2+w_{B_1}\}$; thus every $\mu \notin \mathcal{T}_{Q,w}(K_n)$ has multiplicity at most~$1$.
    \item In case~(III) with $s = 1$, $G = K_{1+b_1}$ has two blocks $A_1$ ($a_1 = 1$) and $B_1$ ($b_1 \ge 1$) with $w_{B_1} - w_{A_1} \ge 1$. Since $\mu \notin \mathcal{T}_{Q,w}(G)$, any $\mu$-eigenvector is constant on $B_1$ and satisfies the $2 \times 2$ quotient system
    \[
      \begin{pmatrix} b_1 + w_{A_1} & b_1 \\ 1 & 2b_1 - 1 + w_{B_1} \end{pmatrix}
      \begin{pmatrix} x_{A_1} \\ x_{B_1} \end{pmatrix} = \mu \begin{pmatrix} x_{A_1} \\ x_{B_1} \end{pmatrix}.
    \]
    Because both off-diagonal entries $b_1, 1$ are strictly positive, the discriminant is $((b_1 - 1 + w_{B_1} - w_{A_1}))^2 + 4b_1 > 0$, so its two eigenvalues are distinct and $\mult_{Q_w(G)}(\mu) \le 1$.
  \end{itemize}

  \emph{General case (case~(I) with $s \ge 1$, or cases~(II)--(IV) with $s \ge 2$).} Because $\mu \notin \mathcal{T}_{Q,w}(G)$ eliminates all twin-difference eigenvectors within each block of size $\ge 2$, every $\mu$-eigenvector $\mathbf{x}$ of $Q_w(G)$ is constant on each block, taking values $x_{A_i}$ on $A_i$ ($i \ge i_0$) and $x_{B_j}$ on $B_j$ ($j \ge 1$). Set $P_0 := 0$, and in cases~(II) and~(IV) set $a_1 := 0, x_{A_1} := 0, P_1 := 0$ so that the recurrence starts cleanly at $B_1$ (with $(\mu - \gamma_{B_1}) x_{B_1} = S_1$ and $S_2 = S_1 - b_1 x_{B_1}$) before proceeding to $A_2, B_2, \dots, A_s, B_s$. Define partial sums
  \[
    P_k = \sum_{i=i_0}^{k} a_i x_{A_i} \quad (1 \le k \le s), \qquad
    S_k = \sum_{j=k}^{s} b_j x_{B_j} \quad (1 \le k \le s), \qquad S_{s+1} = 0.
  \]
  For each existing block, $Q_w(G)\mathbf{x} = \mu \mathbf{x}$ reads
  \begin{align}
    (\mu - \gamma_{A_k}) x_{A_k} &= S_k, & \gamma_{A_k} &:= d_{A_k} + w_{A_k} && (i_0 \le k \le s), \label{eq:thresh_A}\\
    (\mu - \gamma_{B_k}) x_{B_k} &= P_k + S_1, & \gamma_{B_k} &:= d_{B_k} - 1 + w_{B_k} && (1 \le k \le s), \label{eq:thresh_B}
  \end{align}
  with $P_k = P_{k-1} + a_k x_{A_k}$, $S_{k+1} = S_k - b_k x_{B_k}$, and boundary condition $S_{s+1} = 0$.

  We verify the ordering of the shifts $\gamma_{A_k}$ and $\gamma_{B_k}$. Since $d_{A_k} = \sum_{j=k}^s b_j$ is strictly decreasing in $k$ and $w_{A_k}$ is non-increasing in $k$, we have $\gamma_{A_s} < \dots < \gamma_{A_{i_0}}$. Since $d_{B_k} - 1 = \sum_{i=i_0}^k a_i + \sum_{j=1}^s b_j - 2$ is strictly increasing in $k$ and $w_{B_k}$ is non-decreasing in $k$, we have $\gamma_{B_1} < \dots < \gamma_{B_s}$. Comparing $\gamma_{A_{i_0}}$ and $\gamma_{B_1}$:
  \begin{itemize}
    \item In case~(I) ($a_1 \ge 2$): $\gamma_{B_1} - \gamma_{A_1} = (a_1 - 2) + (w_{B_1} - w_{A_1}) \ge 0$, with equality iff $a_1 = 2$ and $w_{B_1} = w_{A_1}$. If equality holds, then since $a_1 = 2 \ge 2$, $A_1$ is an independent $(G,w)$-twin class of size $\ge 2$, so $\gamma_{A_1} = \gamma_{B_1} \in \mathcal{T}_{Q,w}(G)$, which is excluded by $\mu \notin \mathcal{T}_{Q,w}(G)$.
    \item In case~(II) ($a_1 = 0, b_1 \ge 2, s \ge 2$): $\gamma_{B_1} - \gamma_{A_2} = (b_1 - 2) + (w_{B_1} - w_{A_2}) \ge 0$, with equality iff $b_1 = 2$ and $w_{B_1} = w_{A_2}$. If equality holds, then since $b_1 = 2 \ge 2$, $B_1$ is a clique $(G,w)$-twin class of size $\ge 2$, so $\gamma_{B_1} = \gamma_{A_2} \in \mathcal{T}_{Q,w}(G)$, which is excluded by $\mu \notin \mathcal{T}_{Q,w}(G)$.
    \item In case~(III) ($a_1 = 1, w_{B_1} - w_{A_1} \ge 1$): $d_{B_1} - 1 = d_{A_1} - 1$, so $\gamma_{B_1} - \gamma_{A_1} = (w_{B_1} - w_{A_1}) - 1 \ge 0$, with equality iff $w_{B_1} - w_{A_1} = 1$. Moreover, if $\mu = \gamma_{A_1} = \gamma_{B_1}$, then we must have $b_1 = 1$ (for if $b_1 \ge 2$, then $B_1$ would be a clique $(G,w)$-twin class of size $\ge 2$ with local eigenvalue $\gamma_{B_1} \in \mathcal{T}_{Q,w}(G)$, contradicting $\mu \notin \mathcal{T}_{Q,w}(G)$).
    \item In case~(IV) ($a_1 = 0, b_1 = 1, s \ge 2, w_{B_1} - w_{A_2} \ge 1$): $d_{B_1} = \sum_{j=2}^s b_j = d_{A_2}$, so $\gamma_{B_1} - \gamma_{A_2} = (d_{B_1} - 1 + w_{B_1}) - (d_{A_2} + w_{A_2}) = (w_{B_1} - w_{A_2}) - 1 \ge 0$, with equality iff $w_{B_1} - w_{A_2} = 1$. Moreover, if $\mu = \gamma_{A_2} = \gamma_{B_1}$, then we must have $a_2 = 1$ (for if $a_2 \ge 2$, then $A_2$ would be an independent $(G,w)$-twin class of size $\ge 2$ with local eigenvalue $\gamma_{A_2} \in \mathcal{T}_{Q,w}(G)$, contradicting $\mu \notin \mathcal{T}_{Q,w}(G)$).
  \end{itemize}
  Likewise, if $\mu = \gamma_{A_r}$ then $a_r = 1$ (else $A_r$ yields $\gamma_{A_r} \in \mathcal{T}_{Q,w}(G)$), and if $\mu = \gamma_{B_r}$ then $b_r = 1$ (else $B_r$ yields $\gamma_{B_r} \in \mathcal{T}_{Q,w}(G)$). Because $\gamma_{A_s} < \dots < \gamma_{A_{i_0}} \le \gamma_{B_1} < \dots < \gamma_{B_s}$ and $\mu \notin \mathcal{T}_{Q,w}(G)$, there are five exhaustive cases:
  \begin{enumerate}
    \item \emph{No shift equals $\mu$.} Then~\eqref{eq:thresh_A}--\eqref{eq:thresh_B} determine all block values uniquely as linear functions of the single scalar $S_1$, so $\mult_{Q_w(G)}(\mu) \le 1$.
    \item \emph{$\mu = \gamma_{A_r}$ uniquely (so $a_r = 1$, $i_0 \le r \le s$).} Equation~\eqref{eq:thresh_A} at $k=r$ forces $S_r = 0$, and the formal sequence is uniquely determined linearly by $(S_1, x_{A_r})$ subject to the boundary condition $S_{s+1} = 0$. Testing $(S_1, x_{A_r}) = (0, 1)$ sets all entries before $A_r$ to $0$ (so $P_{r-1} = 0, S_r = 0$) and gives $P_r = 1 > 0$. Since $\mu = \gamma_{A_r} < \gamma_{B_k}$ ($k \ge r$) and $\gamma_{A_k} < \mu$ ($k > r$), we have $\mu - \gamma_{B_k} < 0$ ($k \ge r$) and $\mu - \gamma_{A_k} > 0$ ($k > r$). Thus $x_{B_r} = P_r/(\mu - \gamma_{B_r}) < 0$ and $S_{r+1} = -b_r x_{B_r} > 0$. By induction for $k = r+1,\dots,s$:
    \begin{align*}
      x_{A_k} &= \frac{S_k}{\mu - \gamma_{A_k}} > 0 \;\implies\; P_k = P_{k-1} + a_k x_{A_k} > 0 \\
      &\implies\; x_{B_k} = \frac{P_k}{\mu - \gamma_{B_k}} < 0 \;\implies\; S_{k+1} = S_k - b_k x_{B_k} > S_k > 0.
    \end{align*}
    Hence $S_{s+1} > 0$ for $(0, 1)$, so $S_{s+1} = 0$ is a non-trivial linear constraint on $(S_1, x_{A_r})$, leaving $\mult_{Q_w(G)}(\mu) \le 2 - 1 = 1$.
    \item \emph{$\mu = \gamma_{B_r}$ uniquely (so $b_r = 1$, $1 \le r \le s$).} The formal sequence is determined linearly by $(S_1, x_{B_r})$ subject to $S_{s+1} = 0$. Testing $(S_1, x_{B_r}) = (0, -1)$ sets all entries up to $A_r$ to $0$ (so $P_r = 0, S_r = 0$) and yields $S_{r+1} = -b_r(-1) = 1 > 0$. If $r = s$, $S_{s+1} = 1 > 0$ immediately; if $r < s$, then for every $k > r$ we have $\gamma_{A_k} < \mu < \gamma_{B_k}$, so the identical induction gives $S_{s+1} > \dots > S_{r+1} = 1 > 0$. Thus $S_{s+1} = 0$ eliminates one degree of freedom and $\mult_{Q_w(G)}(\mu) \le 1$.
    \item \emph{Case~(III) with $\mu = \gamma_{A_1} = \gamma_{B_1}$ (forcing $a_1 = b_1 = 1$ and $w_{B_1} - w_{A_1} = 1$).} Equation~\eqref{eq:thresh_A} at $k=1$ forces $S_1 = 0$, and equation~\eqref{eq:thresh_B} at $k=1$ forces $P_1 + S_1 = x_{A_1} = 0$. Thus $x_{B_1}$ is the sole remaining free parameter (determining $S_2 = -x_{B_1}$ and all subsequent entries uniquely), so $\mult_{Q_w(G)}(\mu) \le 1$.
    \item \emph{Case~(IV) with $\mu = \gamma_{A_2} = \gamma_{B_1}$ (forcing $b_1 = a_2 = 1$ and $w_{B_1} - w_{A_2} = 1$).} Equation~\eqref{eq:thresh_B} at $k=1$ (with $P_1 = 0$) forces $S_1 = 0$, whence $S_2 = S_1 - x_{B_1} = -x_{B_1}$; equation~\eqref{eq:thresh_A} at $k=2$ then forces $S_2 = 0$, so $x_{B_1} = 0$. Consequently, $x_{A_2}$ is the sole remaining free parameter (determining $P_2 = x_{A_2}$ and all subsequent entries $x_{B_2}, x_{A_3}, \dots, x_{B_s}$ uniquely, since $\mu < \gamma_{B_k}$ for all $k \ge 2$ and $\gamma_{A_k} < \mu$ for all $k \ge 3$), giving $\mult_{Q_w(G)}(\mu) \le 1$. \qedhere
  \end{enumerate}
\end{proof}

\begin{corollary}\label{cor:repaired}
  For every connected threshold graph $G$ of order $n \ge 2$ and every $\theta \notin \mathcal{T}^Q_{\max}(G)$, $\mult_{D^Q(G)}(\theta) \le 2 = \nabla(G) + 1$, and this bound is sharp for $n = 3k+5$ ($k \ge 1$).
\end{corollary}

\begin{proof}
  Immediate from Theorem~\ref{thm:Qsimple} ($w \equiv 0$) and Lemma~\ref{lem:rank1}(2)--(3); sharpness follows from Theorem~\ref{thm:disproof}.
\end{proof}

\subsection{Cotree Decomposition, Inheritance Lemmas, and Universal Linear Upper Bounds}

We now isolate the three structural lemmas underpinning the canonical cotree induction for general connected cographs.

\begin{lemma}[Canonical Cotree Shell--Core Decomposition]\label{lem:shell_core}
  Let $G$ be a connected cograph with $\nabla(G) = k \ge 2$. Its canonical cotree $\mathcal{T}(G)$ determines a shell--core partition $V(G) = V(T) \sqcup V(G^*)$ into an outer \emph{shell} $T$ (possibly empty) and a \emph{branching core} $G^*$ ($|G^*| \ge 2$) such that:
  \begin{enumerate}
    \item $\nabla(G^*) = k$, and $G^*$ has one of two pure branching forms:
    \begin{itemize}
      \item \emph{Core Type U:} $G^* = G_1 \cup \dots \cup G_r$ ($r \ge 2$), where each $G_i$ is a nontrivial connected cograph ($|G_i| \ge 2$), $\sum_{i=1}^r \nabla(G_i) = k$, and $1 \le \nabla(G_i) \le k - 1$ for every $i$;
      \item \emph{Core Type J:} $G^* = H_1 \vee \dots \vee H_r$ ($r \ge 2$), where each $H_i$ is a nontrivial disconnected cograph ($|H_i| \ge 2$), $\sum_{i=1}^r \nabla(H_i) = k$, and $1 \le \nabla(H_i) \le k - 1$ for every $i$.
    \end{itemize}
    \item Whenever $T \ne \emptyset$, $T$ is a connected threshold graph with alternating independent layers $A_1,\dots,A_s$ ($|A_i|=a_i$, $a_1 \ge 0$, $a_2,\dots,a_s \ge 1$) and clique layers $B_1,\dots,B_s$ ($|B_i|=b_i \ge 1$, $s \ge 1$) satisfying $A_i \sim B_j \iff i \le j$, such that every vertex of $W_T := B_1 \cup \dots \cup B_s$ is adjacent to all of $G^*$ and every vertex of $I_T := A_1 \cup \dots \cup A_s$ is non-adjacent to all of $G^*$.
  \end{enumerate}
\end{lemma}

\begin{proof}
  Let $\mathcal{T}(G)$ be the canonical cotree of $G$: its root is labeled $\vee$ (as $G$ is connected), internal node labels alternate strictly between $\vee$ and $\cup$ along every root-to-leaf path, every internal node has at least two children, and leaves represent vertices of $G$.

  Start at the root $u_0$ of $\mathcal{T}(G)$ and descend as long as the current internal node $u$ has at most one nontrivial child (i.e., at most one child subtree with $\ge 2$ vertices):
  \begin{itemize}
    \item If $u$ had $0$ nontrivial children, its subtree would be complete or edgeless, so by Lemma~\ref{lem:cotree_nabla} every ancestor along the path (having only $1$ nontrivial child) would also have $\nabla = 1$, contradicting $\nabla(G) = k \ge 2$.
    \item Since $\mathcal{T}(G)$ is finite, this descent terminates at a unique first internal node $u^*$ having $r \ge 2$ nontrivial children.
  \end{itemize}
  At every strictly visited ancestor above $u^*$, there is exactly one nontrivial child and hence (since every internal node has $\ge 2$ children) at least one leaf child; we collect those leaf children into a clique layer $B_j$ at each $\vee$-node and an independent layer $A_j$ at each $\cup$-node. Crucially, at the terminating node $u^*$ itself, if $u^*$ also possesses leaf children (e.g., $G = K_b \vee (H_1 \vee H_2)$), we \emph{also} peel all leaf children of $u^*$ into the innermost shell layer ($B_1$ if $u^*$ is $\vee$, or $A_1$ if $u^*$ is $\cup$). Define $G^*$ as the induced subgraph on the $r \ge 2$ nontrivial children of $u^*$, and set $T := G - G^*$.

  \textbf{(1)} At every peeled node (including $u^*$), deleting leaf children removes only universal vertices (at $\vee$-nodes) or isolated vertices (at $\cup$-nodes), which preserves $\nabla$ by Lemma~\ref{lem:cotree_nabla}(2)--(3). Hence $\nabla(G^*) = \nabla(G) = k$. Because all leaf children of $u^*$ were peeled into $T$, $G^*$ is purely the disjoint union of $r \ge 2$ nontrivial connected cographs (if $u^*$ is $\cup$, Core Type U) or the join of $r \ge 2$ nontrivial disconnected cographs (if $u^*$ is $\vee$, Core Type J). By Lemma~\ref{lem:cotree_nabla}, $\nabla(G^*) = \sum_{i=1}^r \nabla(G_i) = k$ (Type U) or $\sum_{i=1}^r \nabla(H_i) = k$ (Type J); since $r \ge 2$ and every nontrivial child has $\nabla \ge 1$, each child satisfies $1 \le \nabla \le k - 1$.

  \textbf{(2)} Suppose $T \ne \emptyset$. Because the root $u_0$ is labeled $\vee$, the outermost peeled layer is a non-empty clique layer $B_s$ ($b_s \ge 1$) of universal vertices of $G$ (if $u^* \ne u_0$, then $u_0$ has only one nontrivial child and thus $\ge 1$ leaf child; if $u^* = u_0$, then $T \ne \emptyset$ means $u_0$ has $\ge 1$ leaf child). Indexing peeled layers from innermost ($1$) to outermost ($s$) gives alternating layers $A_1, B_1, \dots, A_s, B_s$ with $b_1,\dots,b_s \ge 1$, $a_2,\dots,a_s \ge 1$, and $a_1 \ge 0$ ($a_1 \ge 1$ iff the innermost non-empty peeled layer comes from a $\cup$-node, i.e., either $u^*$ is $\cup$ with leaf children, or $u^*$ is $\vee$ with $u^* \ne u_0$ and no leaf children). By the cotree adjacency rule, each $A_i$ is adjacent to $B_j$ iff $B_j$ lies strictly above $A_i$ (i.e., $i \le j$), every $B_j$ lies above $G^*$ (so $W_T$ is completely adjacent to $G^*$), and every $A_i$ lies at a $\cup$-node above or at $G^*$ (so $I_T$ is completely non-adjacent to $G^*$). Since $B_s \ne \emptyset$ is universal in $T$, $T$ is a connected threshold graph.
\end{proof}

\begin{lemma}[Weighted Twin and Local Eigenvalue Inheritance on Modules]\label{lem:inherit}
  Let $G$ be a cograph with integer vicinal-monotone weight $w : V(G) \to \mathbb{Z}$, and let $H \subseteq G$ be a \emph{module} of $G$ (meaning every $u \in V(H)$ has the identical external neighborhood $E_H := N_G(u) \setminus V(H)$ of size $c := |E_H|$). Define $w_H : V(H) \to \mathbb{Z}$ by $w_H(u) := w(u) + c$. Then:
  \begin{enumerate}
    \item For all $u, v \in V(H)$, $u \preceq_H v \iff u \preceq_G v$, and $u, v$ are $(H, w_H)$-twins iff they are $(G, w)$-twins.
    \item $w_H$ is integer-valued and vicinal-monotone on $H$, and $\mathcal{T}_{Q, w_H}(H) \subseteq \mathcal{T}_{Q, w}(G)$.
  \end{enumerate}
  In particular, under the decomposition of Lemma~\ref{lem:shell_core}, property~(2) holds for each of the following restrictions:
  \begin{itemize}
    \item $G \to G^*$ with $w^* := w|_{G^*} + |W_T|\mathbf{1}_{G^*}$;
    \item $G^* \to G_i$ (Core Type U) with $w^*|_{G_i}$;
    \item $G^* \to H_i$ (Core Type J) with $w_i := w^*|_{H_i} + (|G^*| - |H_i|)\mathbf{1}_{H_i}$, and $H_i \to C_j$ (for each connected component $C_j$ of $H_i$) with $w_i|_{C_j}$.
  \end{itemize}
\end{lemma}

\begin{proof}
  \textbf{(1)} For any $u, v \in V(H)$, since $N_G(u) = N_H(u) \sqcup E_H$ and $N_G[v] = N_H[v] \sqcup E_H$, we have $N_G(u) \subseteq N_G[v] \iff N_H(u) \subseteq N_H[v]$, and $N_G(u)\setminus\{v\} = N_G(v)\setminus\{u\} \iff N_H(u)\setminus\{v\} = N_H(v)\setminus\{u\}$. Because $w_H(u) - w_H(v) = w(u) - w(v)$, $(H, w_H)$-twins coincide with $(G, w)$-twins inside $V(H)$.

  \textbf{(2)} Vicinal monotonicity of $w_H$ on $H$ is immediate from~(1). Every maximal $(H, w_H)$-twin class $C \subseteq V(H)$ of size $|C| \ge 2$ is contained in a maximal $(G, w)$-twin class of size $\ge 2$ of the same type (clique or independent). For any $u \in C$, $d_G(u) = d_H(u) + c$, so $d_H(u) + w_H(u) = d_G(u) + w(u)$ and $d_H(u) - 1 + w_H(u) = d_G(u) - 1 + w(u)$. Thus every local eigenvalue of $(H, w_H)$ lies in $\mathcal{T}_{Q,w}(G)$.

  The three concrete applications follow directly because $G^*$ is a module in $G$ with $E_{G^*} = W_T$; each $G_i$ is a module in $G^*$ with $E_{G_i} = \emptyset$; each $H_i$ is a module in $G^*$ with $E_{H_i} = G^* \setminus H_i$; and each connected component $C_j$ of $H_i$ is a module in $H_i$ with $E_{C_j} = \emptyset$.
\end{proof}

Because the shell $T$ is \emph{not} a module in $G$ (when $I_T \ne \emptyset$ and $W_T \ne \emptyset$), its restriction requires a separate lemma.

\begin{lemma}[Shell Restriction Lemma]\label{lem:shell_restrict}
  In the setting of Lemma~\ref{lem:shell_core}, let $w : V(G) \to \mathbb{Z}$ be an integer vicinal-monotone weight on $G$, suppose $T \ne \emptyset$, and define $w_T : V(T) \to \mathbb{Z}$ by
  \[
    w_T := w|_T + |G^*|\,\mathbf{1}_{W_T}.
  \]
  Then:
  \begin{enumerate}
    \item $T$ is a connected threshold graph with layers $A_1, B_1, \dots, A_s, B_s$ ($a_1 \ge 0$, all other $a_i, b_i \ge 1$), and $w_T$ is constant on each layer and non-decreasing along $A_s \preceq_T \dots \preceq_T A_{i_0} \preceq_T B_1 \preceq_T \dots \preceq_T B_s$.
    \item If $a_1 = 1$, then $w_T(B_1) - w_T(A_1) \ge |G^*| \ge 2$, so $A_1$ and $B_1$ are \emph{not} $(T, w_T)$-twins (placing $(T, w_T)$ in case~(III) of Lemma~\ref{lem:thresh_canon}). Similarly, if $a_1 = 0$, $b_1 = 1$, and $s \ge 2$, then $w_T(B_1) - w_T(A_2) \ge |G^*| \ge 2$, so $B_1$ and $A_2$ are \emph{not} $(T, w_T)$-twins (placing $(T, w_T)$ in case~(IV) of Lemma~\ref{lem:thresh_canon}).
    \item $\mathcal{T}_{Q, w_T}(T) \subseteq \mathcal{T}_{Q, w}(G)$.
    \item Consequently, for every $\mu \notin \mathcal{T}_{Q, w}(G)$, $\dim\ker(Q_{w_T}(T) - \mu I) \le 1$.
  \end{enumerate}
\end{lemma}

\begin{proof}
  \textbf{(1)} Connectivity and the block structure of $T$ are given by Lemma~\ref{lem:shell_core}(2). In $G$, all vertices within a single layer $A_i$ or $B_j$ are twins, so $w$ (and hence $w_T$) is constant on each layer. Moreover, in $G$ we have $A_s \preceq_G \dots \preceq_G A_{i_0} \preceq_G B_1 \preceq_G \dots \preceq_G B_s$ (every $A_i$ is non-adjacent to $G^*$ and adjacent only to $B_i \cup \dots \cup B_s$, while every $B_j$ is adjacent to $G^*$ and to $A_1 \cup \dots \cup A_j$ and all other $B_{j'}$). Hence $w(A_s) \le \dots \le w(A_{i_0}) \le w(B_1) \le \dots \le w(B_s)$. Adding $0$ on $I_T$ and $+|G^*| \ge 2$ on $W_T$ strictly preserves this chain:
  \[
    w_T(A_s) \le \dots \le w_T(A_{i_0}) < w_T(B_1) \le \dots \le w_T(B_s) \qquad (\text{with the } A\text{-terms vacuous when } I_T = \emptyset).
  \]

  \textbf{(2)} If $a_1 = 1$, then $A_1 \preceq_G B_1$ implies $w(A_1) \le w(B_1)$, so $w_T(B_1) - w_T(A_1) = (w(B_1) - w(A_1)) + |G^*| \ge |G^*| \ge 2 \ge 1$; even though $A_1$ and $B_1$ become unweighted clique twins in $T$ upon deleting $G^*$, their distinct weights prevent them from being $(T, w_T)$-twins. Likewise, if $a_1 = 0$, $b_1 = 1$, and $s \ge 2$, then $A_2 \preceq_G B_1$ implies $w(A_2) \le w(B_1)$, so $w_T(B_1) - w_T(A_2) = (w(B_1) - w(A_2)) + |G^*| \ge |G^*| \ge 2 \ge 1$; even though $B_1$ and $A_2$ become unweighted independent twins in $T$, their distinct weights prevent them from being $(T, w_T)$-twins.

  \textbf{(3)} Let $C$ be any maximal $(T, w_T)$-twin class of size $|C| \ge 2$. No two distinct layers $A_i, A_{i'}$ ($i < i'$) can be twins in $T$ (since $B_i$ is adjacent to $A_i$ and not to $A_{i'}$), and no two distinct layers $B_j, B_{j'}$ ($j < j'$) can be twins in $T$ (since $A_{j'}$ is adjacent to $B_{j'}$ and not to $B_j$, as $a_{j'} \ge 1$ for $j' \ge 2$). Furthermore, no $A_i$ and $B_j$ can be $(T, w_T)$-twins: by~(1), $w_T(A_i) < w_T(B_j)$ holds for \emph{every} $i \in \{i_0,\dots,s\}$ and $j \in \{1,\dots,s\}$ (since $w_T(B_j) - w_T(A_i) \ge |G^*| \ge 2 > 0$), ruling out both unweighted twin pairs $(A_1, B_1)$ (when $a_1 = 1$) and $(A_2, B_1)$ (when $a_1 = 0, b_1 = 1, s \ge 2$) as well as all other $(A_i, B_j)$ pairs. Hence $C$ must be a single layer $A_i$ or $B_j$ (of size $\ge 2$), which is already contained in a maximal $(G, w)$-twin class of size $\ge 2$ of the same type. For $u \in A_i$, $d_T(u) = d_G(u)$ and $w_T(u) = w(u)$, so $d_T(u) + w_T(u) = d_G(u) + w(u)$; for $u \in B_j$, $d_T(u) = d_G(u) - |G^*|$ and $w_T(u) = w(u) + |G^*|$, so $d_T(u) - 1 + w_T(u) = d_G(u) - 1 + w(u)$. Thus $\mathcal{T}_{Q,w_T}(T) \subseteq \mathcal{T}_{Q,w}(G)$.

  \textbf{(4)} By~(3), $\mu \notin \mathcal{T}_{Q,w}(G)$ implies $\mu \notin \mathcal{T}_{Q,w_T}(T)$. If $|T| = 1$ ($T = K_1$), then $\dim\ker(Q_{w_T}(T) - \mu I) \le 1$ trivially; if $|T| \ge 2$, then~(1)--(2) and Theorem~\ref{thm:Qsimple} give $\dim\ker(Q_{w_T}(T) - \mu I) \le 1$.
\end{proof}

\begin{theorem}[Universal Linear Upper Bounds for All Connected Cographs]\label{thm:cograph_upper}
  Let $G$ be any connected cograph on $n \ge 2$ vertices. Then:
  \begin{enumerate}
    \item for every $\mu \notin \mathcal{T}_{Q,\max}(G)$, $\mult_{Q(G)}(\mu) \le 5\nabla(G) - 4$;
    \item for every $\theta \notin \mathcal{T}^Q_{\max}(G)$, $\mult_{D^Q(G)}(\theta) \le 5\nabla(G) - 3$.
  \end{enumerate}
\end{theorem}

\begin{proof}
  We prove by strong induction on $k = \nabla(G)$ the stronger weighted statement: for every connected cograph $G$ on $n \ge 2$ vertices with $\nabla(G) = k$ and every integer vicinal-monotone vertex weight $w : V(G) \to \mathbb{Z}$, every $\mu \notin \mathcal{T}_{Q,w}(G)$ satisfies $\mult_{Q_w(G)}(\mu) \le 5k - 4$.

  \emph{Base case $k = 1$.} Then $G$ is a connected threshold graph, and Theorem~\ref{thm:Qsimple} (via Lemma~\ref{lem:thresh_canon}) gives $\mult_{Q_w(G)}(\mu) \le 1 = 5(1) - 4$.

  \emph{Inductive step ($k \ge 2$).} Apply Lemma~\ref{lem:shell_core} to obtain the shell--core decomposition $V(G) = V(T) \sqcup V(G^*)$ with $\nabla(G^*) = k$. Let $U := \ker(Q_w(G) - \mu I)$ and set
  \[
    U_1 := \Bigl\{\mathbf{x} \in U : \sum_{u \in G^*} x_u = 0\Bigr\},
  \]
  so $\dim U_1 \ge \dim U - 1$.

  \emph{Step 1: Peeling the shell $T$.}
  \begin{itemize}
    \item If $T = \emptyset$, set $U_0 := U_1$; then $\dim U_0 = \dim U_1 \ge \dim U - 1 \ge \dim U - 2$.
    \item If $T \ne \emptyset$, then for every $\mathbf{x} \in U_1$, the cross-term from $G^*$ to $W_T$ vanishes because $\sum_{u \in G^*} x_u = 0$, while $I_T$ has no edges to $G^*$; thus $\mathbf{x}|_T \in \ker(Q_{w_T}(T) - \mu I)$, where $w_T = w|_T + |G^*|\mathbf{1}_{W_T}$. By Lemma~\ref{lem:shell_restrict}(4), $\dim\ker(Q_{w_T}(T) - \mu I) \le 1$. Defining the restriction map $\rho_T : U_1 \to \ker(Q_{w_T}(T) - \mu I)$ by $\rho_T(\mathbf{x}) = \mathbf{x}|_T$ and setting $U_0 := \ker \rho_T = \{\mathbf{x} \in U_1 : \mathbf{x}|_T = \mathbf{0}\}$, rank-nullity gives $\dim U_0 \ge \dim U_1 - 1 \ge \dim U - 2$.
  \end{itemize}
  In both cases, for every $\mathbf{x} \in U_0$, we have $\mathbf{x}|_T = \mathbf{0}$ and $\sum_{u \in G^*} x_u = 0$, so $\mathbf{x} \mapsto \mathbf{x}|_{G^*}$ embeds $U_0$ injectively into $\ker(Q_{w^*}(G^*) - \mu I)$, where $w^* := w|_{G^*} + |W_T|\mathbf{1}_{G^*}$. By Lemma~\ref{lem:inherit}, $w^*$ is integer vicinal-monotone on $G^*$ and $\mu \notin \mathcal{T}_{Q,w^*}(G^*)$.

  \emph{Step 2: Bounding $\dim U_0$ on the branching core $G^*$.}
  \begin{itemize}
    \item \emph{Core Type U} ($G^* = G_1 \cup \dots \cup G_r$, $r \ge 2$, $\sum_{i=1}^r \nabla(G_i) = k$, $1 \le \nabla(G_i) \le k-1$): here $Q_{w^*}(G^*) = \bigoplus_{i=1}^r Q_{w^*|_{G_i}}(G_i)$. By Lemma~\ref{lem:inherit}, each $(G_i, w^*|_{G_i})$ is integer vicinal-monotone with $\mu \notin \mathcal{T}_{Q,w^*|_{G_i}}(G_i)$. Applying the induction hypothesis to each $G_i$ yields
    \[
      \dim U \le 2 + \dim U_0 \le 2 + \sum_{i=1}^r \bigl(5\nabla(G_i) - 4\bigr) = 5k + 2 - 4r \le 5k - 6 < 5k - 4.
    \]
    \item \emph{Core Type J} ($G^* = H_1 \vee \dots \vee H_r$, $r \ge 2$, $\sum_{i=1}^r \nabla(H_i) = k$, $1 \le \nabla(H_i) \le k-1$): for every $\mathbf{x} \in U_0$, $\sum_{u \in G^*} x_u = 0$ implies $\sum_{v \in G^* \setminus H_i} x_v = -\sum_{v \in H_i} x_v$, decoupling the parts $H_1,\dots,H_r$:
    \[
      \mathbf{x}|_{H_i} \in \ker\bigl(Q_{w_i}(H_i) - J_{|H_i|} - \mu I\bigr), \qquad w_i := w^*|_{H_i} + (|G^*| - |H_i|)\mathbf{1}_{H_i}.
    \]
    By Lemma~\ref{lem:inherit}, each $(H_i, w_i)$ is integer vicinal-monotone with $\mu \notin \mathcal{T}_{Q,w_i}(H_i)$. Write $H_i = C_1 \cup \dots \cup C_{t_i} \cup \overline{K_{a_i}}$ ($t_i \ge 0$ nontrivial connected components, $a_i \ge 0$ isolated vertices), and set $d_i := \dim\ker(Q_{w_i}(H_i) - J_{|H_i|} - \mu I)$. We prove $d_i \le 5\nabla(H_i) - 3$:
    \begin{itemize}
      \item \emph{Subcase $t_i = 0$ ($H_i = \overline{K_{a_i}}$, $a_i \ge 2$, $\nabla(H_i) = 1$):} all vertices of $H_i$ are $(H_i, w_i)$-twins with local eigenvalue $w_i(H_i) \ne \mu$, so $Q_{w_i}(H_i) - \mu I = (w_i(H_i) - \mu)I$ is invertible ($\nullity = 0$). By Lemma~\ref{lem:rank1}(1) applied to the rank-one perturbation $-J_{|H_i|}$, $d_i \le 0 + 1 = 1 < 5(1) - 3 = 2$.
      \item \emph{Subcase $t_i \ge 2$:} by Lemma~\ref{lem:cotree_nabla}(2), $\sum_{j=1}^{t_i} \nabla(C_j) = \nabla(H_i) \le k - 1$. By Lemma~\ref{lem:inherit}, each $(C_j, w_i|_{C_j})$ satisfies the induction hypothesis, giving $\nullity(Q_{w_i}(C_j) - \mu I) \le 5\nabla(C_j) - 4$. On $\overline{K_{a_i}}$, $\nullity(Q_{w_i}(\overline{K_{a_i}}) - \mu I)$ is $0$ if $a_i \ge 2$ (since $w_i(\overline{K_{a_i}}) \in \mathcal{T}_{Q,w_i}(H_i)$) and at most $1$ if $a_i = 1$. By Lemma~\ref{lem:rank1}(1) applied to $-J_{|H_i|}$,
      \[
        d_i \le 1 + \mathbf{1}_{a_i = 1} + \sum_{j=1}^{t_i} \bigl(5\nabla(C_j) - 4\bigr) \le 2 + 5\nabla(H_i) - 4t_i \le 5\nabla(H_i) - 6 < 5\nabla(H_i) - 3.
      \]
      \item \emph{Subcase $t_i = 1$ ($H_i = C_1 \cup \overline{K_{a_i}}$, $a_i \ge 1$, $\nabla(C_1) = \nabla(H_i) \le k-1$):} by Lemma~\ref{lem:inherit} and the induction hypothesis, $\nullity(M) \le 5\nabla(H_i) - 4$ where $M := Q_{w_i}(C_1) - \mu I$. If $a_i \ge 2$, then $\nullity(Q_{w_i}(\overline{K_{a_i}}) - \mu I) = 0$ and Lemma~\ref{lem:rank1}(1) directly gives $d_i \le \nullity(M) + 1 \le 5\nabla(H_i) - 3$. If $a_i = 1$ (a single isolated vertex $v_0$), set $\beta := w_i(v_0) - \mu - 1$ and write $Q_{w_i}(H_i) - J_{|H_i|} - \mu I$ in block form over $V(C_1) \sqcup \{v_0\}$:
      \[
        N := \begin{pmatrix} M - J & -\mathbf{1} \\ -\mathbf{1}^T & \beta \end{pmatrix}.
      \]
      \begin{itemize}
        \item If $\beta \ne 0$, eliminating the last row/column via the Schur complement yields
        \[
          d_i = \nullity(N) = \nullity\Bigl(M - \Bigl(1 + \frac{1}{\beta}\Bigr)J\Bigr) \le \nullity(M) + 1 \le 5\nabla(H_i) - 3
        \]
        by Lemma~\ref{lem:rank1}(1).
        \item If $\beta = 0$, then $(\mathbf{x}, y)^T \in \ker N$ ($\mathbf{x} \in \mathbb{R}^{|C_1|}, y \in \mathbb{R}$) satisfies $-\mathbf{1}^T \mathbf{x} = 0$ and $(M - J)\mathbf{x} - y\mathbf{1} = \mathbf{0}$, i.e., $\mathbf{1}^T \mathbf{x} = 0$ and $M\mathbf{x} = y\mathbf{1}$. Define the linear functional $L : \ker N \to \mathbb{R}$ by $L(\mathbf{x}, y) := y$. Then $\dim\im L \le 1$, and every $(\mathbf{x}, 0) \in \ker L$ satisfies $M\mathbf{x} = \mathbf{0}$, so $\dim\ker L \le \nullity(M)$. By rank-nullity, $d_i = \dim\ker N = \dim\ker L + \dim\im L \le \nullity(M) + 1 \le 5\nabla(H_i) - 3$.
      \end{itemize}
    \end{itemize}
    Summing $d_i \le 5\nabla(H_i) - 3$ over all $r \ge 2$ parts yields
    \[
      \dim U \le 2 + \dim U_0 \le 2 + \sum_{i=1}^r \bigl(5\nabla(H_i) - 3\bigr) = 5k + 2 - 3r \le 5k - 4.
    \]
  \end{itemize}
  This completes the induction and proves~(1) (setting $w \equiv 0$). Part~(2) follows immediately from~(1) and Lemma~\ref{lem:rank1}(2)--(3).
\end{proof}

\subsection{Linear Violation of $\nabla(G)$ and $\nabla(G)+C$ on General Connected Cographs}

\begin{theorem}[Linearly Unbounded Separation on Connected Cographs]\label{thm:cograph_lower}
  Let $B := K_{36} \vee (3 H_8)$ be the $60$-vertex connected cograph obtained by joining $K_{36}$ to three disjoint copies of the $8$-vertex antiregular graph $H_8 = C(1^8)$, and for every integer $s \ge 1$ let
  \[
    F_s := K_1 \vee (s \cdot B)
  \]
  be the connected cograph of order $n_s = 60s + 1$. Then:
  \begin{enumerate}
    \item $\nabla(B) = 3$, $40 \notin \mathcal{T}_{Q,\max}(B)$, and $\mult_{Q(B)}(40) = 4 > 3 = \nabla(B)$;
    \item for every $s \ge 1$, $\nabla(F_s) = 3s$, $41 \notin \mathcal{T}_{Q,\max}(F_s)$, $\theta_s := 120s - 43 \notin \mathcal{T}^Q_{\max}(F_s)$, and
    \[
      \mult_{Q(F_s)}(41) \ge 4s - 1 = \frac{4}{3}\nabla(F_s) - 1, \qquad
      \mult_{D^Q(F_s)}(\theta_s) \ge 4s - 1 = \frac{4}{3}\nabla(F_s) - 1.
    \]
  \end{enumerate}
  In particular, $F_3$ ($n = 181$, $\nabla = 9$) has $\mult_{Q(F_3)}(41) \ge 11$ and $\mult_{D^Q(F_3)}(317) \ge 11 = \nabla(F_3) + 2$, refuting both $\mult_Q \le \nabla$ and $\mult_{D^Q} \le \nabla + 1$ (and indeed any $\nabla + C$ bound) on general connected cographs.
\end{theorem}

\begin{proof}
  \textbf{(1)} By Lemma~\ref{lem:cotree_nabla}, $\nabla(B) = 3\nabla(H_8) = 3$. The maximal twin classes of size $\ge 2$ in $B$ are $K_{36}$ (degree $59$, local $Q$-eigenvalue $58$) and the $K_2$ twin class $\{v_1, v_2\}$ in each $H_8$ (degree $4 + 36 = 40$, local $Q$-eigenvalue $39$), so $\mathcal{T}_{Q,\max}(B) = \{58, 39\}$ and $40 \notin \mathcal{T}_{Q,\max}(B)$.

  From Theorem~\ref{thm:disproof} ($k=1$), $H_8$ has $\mathbf{x}^{(1)} = (-1,-1,1,1,0,0,0,0)^T$ with $\mathbf{1}_8^T \mathbf{x}^{(1)} = 0$ and $Q(H_8)\mathbf{x}^{(1)} = 4\mathbf{x}^{(1)}$, and $\mathbf{x}^{(2)} = (2,2,-2,0,-1,1,-1,1)^T$ with $\mathbf{1}_8^T \mathbf{x}^{(2)} = 2$ and $D^Q(H_8)\mathbf{x}^{(2)} = 8\mathbf{x}^{(2)}$. By~\eqref{eq:rank1} on $H_8$ ($2(8)-4 = 12$),
  \[
    \bigl(12I - Q(H_8) + 2J\bigr)\mathbf{x}^{(2)} = 8\mathbf{x}^{(2)} \;\implies\; \bigl(Q(H_8) - 4I\bigr)\mathbf{x}^{(2)} = 2J\mathbf{x}^{(2)} = 4\mathbf{1}_8.
  \]
  In $B = K_{36} \vee (3H_8)$, placing $\mathbf{x}^{(1)}$ on each of the three copies of $H_8$ (and $0$ elsewhere) gives three pairwise disjointly supported $Q(B)$-eigenvectors $\mathbf{u}^{(1)}, \mathbf{u}^{(2)}, \mathbf{u}^{(3)}$ for $\mu = 4 + 36 = 40$, each orthogonal to $\mathbf{1}_{60}$. Define a fourth vector $\mathbf{u}^{(4)} \in \mathbb{R}^{60}$ equal to $-1$ on $K_{36}$ and $9\mathbf{x}^{(2)}$ on each copy of $H_8$. Then:
  \begin{itemize}
    \item on each copy of $H_8$: $(Q(H_8) + 36I)(9\mathbf{x}^{(2)}) + 36(-1)\mathbf{1}_8 = 9(4\mathbf{x}^{(2)} + 4\mathbf{1}_8) + 36(9\mathbf{x}^{(2)}) - 36\mathbf{1}_8 = 40(9\mathbf{x}^{(2)})$;
    \item on $K_{36}$ (degree $59$): $(59 + 35)(-1) + 3 \cdot 9(\mathbf{1}_8^T \mathbf{x}^{(2)}) = -94 + 54 = -40 = 40(-1)$.
  \end{itemize}
  Since $\mathbf{u}^{(4)} \equiv -1 \ne 0$ on $K_{36}$ while $\mathbf{u}^{(1)}, \mathbf{u}^{(2)}, \mathbf{u}^{(3)}$ vanish on $K_{36}$ and have pairwise disjoint supports on the three $H_8$ copies, $\mathbf{u}^{(1)},\dots,\mathbf{u}^{(4)}$ are linearly independent, so $\mult_{Q(B)}(40) \ge 4$.

  To prove the exact equality $\mult_{Q(B)}(40) = 4$ purely analytically, let $E_{40} := \ker(Q(B) - 40 I_{60})$. Because $K_{36}$ is a clique twin class of local $Q$-eigenvalue $58 \ne 40$, every $\mathbf{x} \in E_{40}$ is constant on $K_{36}$, taking a common value $\lambda(\mathbf{x}) \in \mathbb{R}$. Define the linear functional $L : E_{40} \to \mathbb{R}$ by $L(\mathbf{x}) := \lambda(\mathbf{x})$, so $\dim\im L \le 1$. If $\mathbf{x} \in \ker L$, then $\mathbf{x}|_{K_{36}} = \mathbf{0}$, so on each copy of $H_8$ the eigenvector equation reduces to $(Q(H_8) + 36I)\mathbf{x}|_{H_8} = 40\mathbf{x}|_{H_8}$, i.e., $\mathbf{x}|_{H_8} \in \ker(Q(H_8) - 4I)$. Since $H_8$ is a connected threshold graph and $4 \notin \mathcal{T}_{Q,\max}(H_8) = \{3\}$, Theorem~\ref{thm:Qsimple} (or Theorem~\ref{thm:CgraphQ}) gives $\dim\ker(Q(H_8) - 4I) = 1$, whence $\dim\ker L \le 3$. By rank-nullity,
  \[
    \mult_{Q(B)}(40) = \dim E_{40} = \dim\ker L + \dim\im L \le 3 + 1 = 4,
  \]
  proving $\mult_{Q(B)}(40) = 4$.

  \textbf{(2)} In $F_s = K_1 \vee (s \cdot B)$, Lemma~\ref{lem:cotree_nabla} gives $\nabla(F_s) = s\nabla(B) = 3s$, and the maximal twin classes of size $\ge 2$ have local $Q$-eigenvalues $59$ and $40$ (noting that for $s = 1$, $F_1 = K_{37} \vee (3H_8)$ merges $K_1$ with $K_{36}$ into a single $K_{37}$ class of degree $60$ and the same local $Q$-eigenvalue $59$), so $41 \notin \mathcal{T}_{Q,\max}(F_s)$ and $\theta_s = 2(60s+1) - 4 - 41 = 120s - 43 \notin \mathcal{T}^Q_{\max}(F_s)$ by Lemma~\ref{lem:rank1}(3). Consider the following $4s - 1$ vectors in $\mathbb{R}^{60s+1}$, all vanishing on $K_1$ and orthogonal to $\mathbf{1}_{60s+1}$:
  \begin{itemize}
    \item for $j \in \{1,\dots,s\}$ and $t \in \{1,2,3\}$, the vector $\mathbf{v}^{(j,t)}$ equal to $\mathbf{u}^{(t)}$ on copy~$j$ of $B$ and $0$ elsewhere ($3s$ vectors);
    \item for $j \in \{2,\dots,s\}$, the difference vector $\mathbf{w}^{(j)}$ equal to $+\mathbf{u}^{(4)}$ on copy~$1$ of $B$, $-\mathbf{u}^{(4)}$ on copy~$j$ of $B$, and $0$ elsewhere ($s - 1$ vectors).
  \end{itemize}
  To verify linear independence, suppose $\sum_{j=1}^s \sum_{t=1}^3 \alpha_{j,t} \mathbf{v}^{(j,t)} + \sum_{j=2}^s \beta_j \mathbf{w}^{(j)} = \mathbf{0}$. Restricting to the $K_{36}$ block of copy~$j$ ($j \in \{2,\dots,s\}$), every $\mathbf{v}^{(j',t)}$ vanishes on $K_{36}$ and the only vector non-zero there is $\mathbf{w}^{(j)}$ (taking value $+1$ on $K_{36}$ of copy~$j$), forcing $\beta_j = 0$ for all $j = 2,\dots,s$. What remains is $\sum_{j,t} \alpha_{j,t} \mathbf{v}^{(j,t)} = \mathbf{0}$; since the $3s$ vectors $\mathbf{v}^{(j,t)}$ have pairwise disjoint supports (each supported on a distinct copy of $H_8$), $\alpha_{j,t} = 0$ for all $j, t$.

  Because every $\mathbf{v} \in \{\mathbf{v}^{(j,t)}, \mathbf{w}^{(j)}\}$ vanishes on $K_1$ and satisfies $\mathbf{1}_{60s+1}^T \mathbf{v} = 0$, the cross-terms between $K_1$ and $s \cdot B$ vanish, giving $Q(F_s)\mathbf{v} = (40 + 1)\mathbf{v} = 41\mathbf{v}$. Since $J\mathbf{v} = \mathbf{0}$, identity~\eqref{eq:rank1} immediately yields $D^Q(F_s)\mathbf{v} = (2n_s - 4 - 41)\mathbf{v} = \theta_s \mathbf{v}$. Hence both $Q(F_s)$ and $D^Q(F_s)$ attain non-local multiplicity at least $4s - 1 = \frac{4}{3}\nabla(F_s) - 1$.
\end{proof}

\begin{remark}[Exact symbolic verification on small orders and reproducibility]\label{rem:enum}
  Remarkably, the linear separation in Theorem~\ref{thm:cograph_lower} is invisible at small orders: using our companion script \texttt{verify\_all.py} (verified under Python~3.11.2, NumPy~2.1.3 for integer arrays, and SymPy~1.13.3 for exact integer matrices and irreducible-factor decomposition \texttt{factor\_list} over $\mathbb{Q}[x]$, with zero floating-point arithmetic), we exhaustively checked all $1{,}170$ non-isomorphic connected cographs of order $2 \le n \le 9$ via canonical cotrees. Table~\ref{tab:cographs} reports the \emph{number of graphs} exhibiting at least one maximal-twin-non-local eigenvalue violating each bound: \emph{zero} graphs violate $\mult_Q \le \nabla$ or $\mult_{D^Q} \le \nabla + 1$, and exactly $4$ graphs (all at $n = 8$, each at a single non-local eigenvalue) violate $\mult_{D^Q} \le \nabla$. As noted in Section~\ref{sec:intro}, these finite-order symbolic checks confirm exact small-order counts and independently cross-check our explicit constructions, while all universal bounds in Sections~\ref{sec:Cgraph_bound}--\ref{sec:cographs} are proved purely analytically.
\end{remark}

\begin{table}[!htbp]
  \centering
  \caption{Exact symbolic verification across all $1{,}170$ non-isomorphic connected cographs of order $2 \le n \le 9$ (counts denote \emph{number of graphs} with at least one maximal-twin-non-local eigenvalue violating the stated inequality).}
  \label{tab:cographs}
  {\small\setlength{\tabcolsep}{5pt}
  \begin{tabular}{ccccc}
    \toprule
    Order $n$ & Cographs & \#\,Graphs ($\mult_{Q} > \nabla$) & \#\,Graphs ($\mult_{D^Q} > \nabla$) & \#\,Graphs ($\mult_{D^Q} > \nabla + 1$) \\
    \midrule
    $2$ & $1$   & $0$ & $0$ & $0$ \\
    $3$ & $2$   & $0$ & $0$ & $0$ \\
    $4$ & $5$   & $0$ & $0$ & $0$ \\
    $5$ & $12$  & $0$ & $0$ & $0$ \\
    $6$ & $33$  & $0$ & $0$ & $0$ \\
    $7$ & $90$  & $0$ & $0$ & $0$ \\
    $8$ & $261$ & $0$ & $4$ & $0$ \\
    $9$ & $766$ & $0$ & $0$ & $0$ \\
    \midrule
    Total & $1{,}170$ & $0$ & $4$ & $0$ \\
    \bottomrule
  \end{tabular}}
\end{table}

\begin{problem}\label{prob:cstar}
  For each integer $k \ge 1$, define the extremal maximal-twin-non-local multiplicities over connected cographs of Dilworth number $k$:
  \begin{align*}
    M_Q(k) &:= \sup_{\substack{G\text{ connected cograph} \\ \nabla(G) = k}} \max_{\mu \notin \mathcal{T}_{Q,\max}(G)} \mult_{Q(G)}(\mu), \\
    M_{D^Q}(k) &:= \sup_{\substack{G\text{ connected cograph} \\ \nabla(G) = k}} \max_{\theta \notin \mathcal{T}^Q_{\max}(G)} \mult_{D^Q(G)}(\theta),
  \end{align*}
  and define the asymptotic growth constants
  \[
    c_Q^* := \limsup_{k \to \infty} \frac{M_Q(k)}{k}, \qquad
    c_{D^Q}^* := \limsup_{k \to \infty} \frac{M_{D^Q}(k)}{k}.
  \]
  By Lemma~\ref{lem:rank1}(2)--(3), $|M_Q(k) - M_{D^Q}(k)| \le 1$ for every $k \ge 1$, whence $c_Q^* = c_{D^Q}^*$. Theorems~\ref{thm:cograph_upper} and~\ref{thm:cograph_lower} establish
  \[
    \frac{4}{3} \le c_Q^* = c_{D^Q}^* \le 5.
  \]
  Determine the exact value of $c_Q^* = c_{D^Q}^*$.
\end{problem}

\begin{thebibliography}{9}

\bibitem{AH2016}
M.~Aouchiche and P.~Hansen,
\emph{On the distance signless Laplacian of a graph},
Linear Multilinear Algebra \textbf{64} (2016), 1113--1123.

\bibitem{Ghorbani2019}
E.~Ghorbani,
\emph{Cographs: Eigenvalues and Dilworth number},
Discrete Math. \textbf{342} (2019), 2797--2803.

\bibitem{LWB2024}
L.~Li, J.~Wang, and M.~Brunetti,
\emph{Seidel matrices, Dilworth number and an eigenvalue-free interval for cographs},
Linear Algebra Appl. \textbf{698} (2024), 56--72.

\bibitem{LL2021}
Z.~Lou and H.~Lin,
\emph{Distance eigenvalues of cographs and their multiplicities},
Linear Algebra Appl. \textbf{608} (2021), 1--12.

\bibitem{MM2025}
S.~Mandal and R.~Mehatari,
\emph{Spectral properties of $\mathcal{C}$-graphs},
Comput. Appl. Math. \textbf{44} (2025), Article No.~136.

\bibitem{MM2026}
S.~Mandal and P.~Manna,
\emph{The distance Laplacian and distance signless Laplacian spectra of $\mathcal{C}$-graphs},
Preprint, \href{https://arxiv.org/abs/2610.08599}{arXiv:2610.08599} (October~6, 2026).

\bibitem{MMT2026}
P.~Manna, S.~Mandal, and F.~C.~Tura,
\emph{The spectral theory of cographs: A survey, organized by graph matrix},
Preprint, \href{https://arxiv.org/abs/2611.06981}{arXiv:2611.06981} (November~4, 2026).

\bibitem{TrureTuring2026}
The Omega Institute,
\emph{Refutation of Mandal--Manna Conjecture~5.5},
\texttt{trureturing} External Open Problems Registry (frozen November~7, 2026),
\href{https://the-omega-institute.github.io/trureturing-mdbook/open-problems.html}{https://the-omega-institute.github.io/trureturing-mdbook/open-problems.html}.

\end{thebibliography}

\end{document}
